% !TEX program = lualatex
% ============================================================
% ملحق إرنفست: تفسير التناسق لمفارقة القرص الدوار ونشوء الهندسة غير الإقليدية
% الترجمة إلى العربية
% ============================================================

\PassOptionsToPackage{unicode=true}{hyperref}
\documentclass[12pt,a4paper]{article}

% ============================================================
% تعريف أمر \babelsublr للأرقام العربية في سلاسل PDF
% ============================================================
\makeatletter
\ifdefined\babelsublr\else
\newcommand\babelsublr[1]{#1}
\fi
\makeatother

% ============================================================
% إعدادات اللغة العربية والخطوط
% ============================================================
\usepackage{polyglossia}
\setmainlanguage{arabic}
\setotherlanguage{english}

% خطوط عربية
\newfontfamily\arabicfont{Amiri}[Script=Arabic, Language=Arabic]
\newfontfamily\arabicfontsf{Amiri}[Script=Arabic, Language=Arabic]
\newfontfamily\arabicfonttt{Amiri}[Script=Arabic, Language=Arabic]

% خطوط لاتينية
\newfontfamily\englishfont{Times New Roman}[Script=Latin, Language=English]
\setmonofont{Courier New}[Scale=0.85]

% ============================================================
% الحزم الأساسية
% ============================================================
\usepackage{amsmath,amssymb,amsthm}
\usepackage{geometry}
\geometry{margin=1in}
\usepackage{hyperref}
\usepackage{booktabs}
\usepackage{graphicx}
\usepackage{tikz}
\usetikzlibrary{arrows.meta, positioning, shapes.geometric, decorations.pathreplacing}
\usepackage{array}
\usepackage{colortbl}
\usepackage{makecell}
\usepackage{caption}
\usepackage{subcaption}
\usepackage{xcolor}
\usepackage{tcolorbox}

\hypersetup{colorlinks=true, linkcolor=blue, citecolor=blue, urlcolor=blue}

% ============================================================
% بيئات نظرية
% ============================================================
\newtheorem{theorem}{نظرية}[section]
\newtheorem{definition}{تعريف}[section]
\newtheorem{corollary}{نتيجة}[section]
\newtheorem{proposition}{افتراض}[section]
\newtheorem{remark}{ملاحظة}[section]
\newtheorem{prediction}{تنبؤ}[section]

% ============================================================
% ماكروت YUCT
% ============================================================
\newcommand{\kappac}{\kappa_c}
\newcommand{\eps}{\varepsilon}
\newcommand{\betaf}{\beta}
\newcommand{\dint}{D\!+\!I\!\cdot\!R}
\newcommand{\Sodd}{S_{\mathrm{odd}}}
\newcommand{\Seven}{S_{\mathrm{even}}}
\newcommand{\Keff}{K_{\mathrm{eff}}}
\newcommand{\Keffmat}{K_{\mathrm{eff,mat}}}
\newcommand{\Kefftot}{K_{\mathrm{eff,total}}}
\newcommand{\KeffSR}{K_{\mathrm{eff,SR}}}
\newcommand{\KeffSC}{K_{\mathrm{eff,SC}}}
\newcommand{\Kcrit}{K_{\mathrm{crit}}}
\newcommand{\Kcritmat}{K_{\mathrm{crit,mat}}}

\begin{document}
	
	\begin{titlepage}
		\begin{center}
			\vspace*{1cm}
			\Huge\textbf{ملحق إرنفست: تفسير التناسق لمفارقة القرص الدوار ونشوء الهندسة غير الإقليدية} \\
			\vspace{0.5cm}
			\Large\textit{من افتراضات الجسم الجاسئ إلى كفاءة التناسق \(K_{\mathrm{eff}}\) وأصل الانحناء}
			\vspace{1.5cm}
			
			\Large أليكسي ف. ياكوشيف\textsuperscript{1} \\
			\vspace{0.3cm}
			\large \textsuperscript{1}أبحاث ياكوشيف، \url{https://yuct.org/} \\
			\vspace{1cm}
			\textbf{DOI: \url{https://doi.org/10.5281/zenodo.18444598}} \\
			\vspace{1cm}
			أبريل 2026 \\ \large الإصدار 4.4 (مع عتبة الانهيار الحرج وتنبؤات معتمدة على المادة)
			\vspace{2cm}
			
			\begin{minipage}{0.9\textwidth}
				\begin{abstract}
					\noindent
					تسلط مفارقة إرنفست (1910) الضوء على استحالة الحفاظ على الهندسة الإقليدية على قرص دوار مع احترام تقلص الطول النسبي. في النسبية العامة، تُحل هذه المفارقة بإدخال هندسة غير إقليدية للراصد الدوار. نقدم هنا تفسيرًا ضمن نظرية ياكوشيف الموحدة للتناسق (YUCT) يستند إلى ثوابت التناسق الشاملة المشتقة في الملحق Ferra1.2. نوضح أن تقلص المحيط الظاهر هو مظهر من مظاهر كفاءة التناسق العالية \(K_{\mathrm{eff}} > 1\) التي تتحقق عبر قاموس موزع مسبقًا للمسافات المتبادلة. الثوابت الأساسية \(\Sodd = 6/5 = 1.2\)، \(\Seven = 4/5 = 0.8\)، مجموعهما \(2\)، نسبتهما \(3/2\)، والعلاقة \(\Sodd \varphi^2 \approx \pi\) (حيث \(\varphi\) هي النسبة الذهبية) توفر العمود الفقري الهندسي. نقدم مفهوم \textbf{مركز التناسق} (محور الدوران) الذي وجوده مشروط بوجود مادة منسقة. عندما تتجاوز السرعة الزاوية عتبة حرجة، تنخفض كفاءة تناسق المادة \(K_{\mathrm{eff,mat}}\) عن قيمة حرجة \(K_{\mathrm{crit}}\)، مما يؤدي إلى فقدان كارثي للتناسق — يتفكك القرص، ويتوقف المحور عن الوجود كخط متميز فيزيائيًا. نشتق تنبؤات كمية للسرعة الزاوية الحرجة بدلالة معاملات مادية قابلة للقياس (معامل يونغ، الكثافة، قوة الشد) ونبين أنه في النهاية \(K_{\mathrm{eff,mat}}=1\) تستعاد الصيغ القياسية لميكانيكا الكسر. تتم مقارنة الإطار مع النظريات البديلة (النسبية العامة، المرونة الكلاسيكية) ويُظهر أنه يقدم تنبؤات جديدة قابلة للاختبار: السرعة الحرجة لقرص فائق التوصيل يجب أن تختلف عن قرص عادي، كنتيجة مباشرة لكفاءة التناسق المعتمدة على المادة. يتم تحديد بروتوكول تجريبي مفصل، وتقديم تقديرات عددية لقرص من موصل فائق عند درجة حرارة عالية.
				\end{abstract}
			\end{minipage}
			\vspace{1cm}
			
			\noindent \textbf{الكلمات المفتاحية:} مفارقة إرنفست، قرص دوار، كفاءة التناسق، YUCT، هندسة غير إقليدية، النسبية العامة، هندسة ناشئة، اختبار الموصلية الفائقة، النسبة الذهبية، ثوابت التناسق الشاملة، ميكانيكا الكسر، الانهيار الحرج.
		\end{center}
	\end{titlepage}
	
	\tableofcontents
	\newpage
	
	\section{مقدمة: مفارقة إرنفست وحلها القياسي}
	\label{sec:intro}
	
	في عام 1910، أشار بول إرنفست إلى تناقض صارخ ضمن إطار النسبية الخاصة عند تطبيقه على قرص جاسئ دوار \cite{Ehrenfest1909}. consider قرص نصف قطره \(R_0\) (كما يقاس في إطار سكونه) يدور بسرعة زاوية ثابتة \(\omega\). وفقًا للنسبية الخاصة، يرى راصد ساكن في المختبر أن كل عنصر مماسي يتقلص طوله بعامل لورنتز \(\gamma = 1/\sqrt{1 - \omega^2 R_0^2/c^2}\). وبالتالي، يصبح المحيط المقاس في المختبر
	\[
	C_{\text{lab}} = 2\pi R_0 \, \gamma^{-1} = 2\pi R_0 \sqrt{1 - \omega^2 R_0^2/c^2},
	\]
	بينما يبقى نصف القطر (العمودي على الحركة) \(R_0\). بالنسبة للمحيط، تفشل العلاقة الإقليدية \(C = 2\pi R\): \(C_{\text{lab}}/(2R_0) = \sqrt{1 - \omega^2 R_0^2/c^2} < \pi\). يُعرف هذا التناقض باسم \textbf{مفارقة إرنفست}.
	
	الحل القياسي ضمن النسبية العامة هو أن الهندسة على القرص الدوار \textbf{غير إقليدية}؛ وهي تتوافق مع فضاء ذي انحناء ثابت، ولا يمكن اعتبار القرص ``جاسئًا'' بمعنى الميكانيكا الكلاسيكية. بالنسبة لراصد مصاحب، فإن المترية هي مترية لانجفان المعروفة، والجزء المكاني منها منحنٍ. وهكذا تقبل النسبية العامة الطبيعة غير الإقليدية كخاصية هندسية أساسية يسببها الدوران.
	
	\subsection{ثوابت التناسق الشاملة كأوليات هندسية}
	
	قبل تقديم تفسير YUCT، نستذكر ثوابت التناسق الشاملة المشتقة في الملحق Ferra1.2 من البنية الهرمية لـ YUCT. تظهر هذه الثوابت من التسلسل الهرمي الكسري للأنظمة المنسقة وتحقق علاقات كسرية دقيقة:
	\begin{align}
		\Sodd &= \frac{6}{5} = 1.2, \qquad \Seven = \frac{4}{5} = 0.8, \\
		\Sodd + \Seven &= 2, \qquad \frac{\Sodd}{\Seven} = \frac{3}{2} = 1.5 = \frac{1}{\betaf},
	\end{align}
	حيث \(\betaf = 2/3 \approx 0.667\) هو أس القياس الشامل. تميز هذه الثوابت المساهمات النهائية للارتباطات أحادية الاتجاه (المستويات الفردية) والمتناوبة (المستويات الزوجية) في أي نظام تناسق هرمي.
	
	هناك علاقة تقريبية رائعة تربط هذه الثوابت بالنسبة الذهبية \(\varphi = (1+\sqrt{5})/2 \approx 1.618034\) وثابت الدائرة \(\pi\):
	\begin{equation}
		\Sodd \cdot \varphi^2 = \frac{6}{5} \cdot \left(\frac{1+\sqrt{5}}{2}\right)^2 = \frac{9+3\sqrt{5}}{5} \approx 3.1416407865,
	\end{equation}
	وهذا يختلف عن \(\pi \approx 3.1415926535\) بنسبة \(4.8\times10^{-5}\) فقط (خطأ نسبي \(1.5\times10^{-5}\)). يشير هذا التقارب القريب إلى وجود رابط هيكلي عميق بين التسلسل الهرمي للتناسق (\(\Sodd\))، والنسبة الذهبية (التشابه الذاتي)، وهندسة الدائرة. في YUCT، يمكن تفسير \(\pi\) على أنه نهاية هذا التعبير عندما تميل كفاءة التناسق \(K_{\mathrm{eff}} \to \infty\) (تماسك تام). بالنسبة لـ \(K_{\mathrm{eff}}\) المحدود، قد ينحرف الثابت الهندسي الفعال عن \(\pi\) بطريقة تخضع لنفس الثوابت الهرمية. هذه الملاحظة ستوجه تفسيرنا للقرص الدوار.
	
	\section{منظور YUCT للتناسق}
	\label{sec:yuct}
	
	تقدم نظرية ياكوشيف الموحدة للتناسق (YUCT) أساسًا أنطولوجيًا مختلفًا. فبدلاً من افتراض الهندسة كأولية، تفترض YUCT أن \textbf{التناسق} — محاذاة الحالات باستخدام قواميس موزعة مسبقًا — هو العملية الأساسية. تنبثق هندسة الزمكان من كفاءة التناسق \(K_{\mathrm{eff}}\)، المُعرَّفة عبر بروتوكول YPSDC (الملحق ب). ينص قانون خطأ القياس الشامل (الملحق ل) على أن أي خطأ نسبي \(\eps\) يحقق
	\[
	\eps = \kappac \alpha K_{\mathrm{eff}}^{-\betaf}, \qquad \betaf = 0.67,\ \kappac = 0.33,
	\]
	ولكن الأهم من ذلك، أن مفهوم ``المسافة'' ذاته هو نتيجة تفاعلات منسقة. توفر الثوابت \(\Sodd\) و \(\Seven\) العمود الفقري الكمي: النسبة \(3/2\) تحكم الانتقال من الطور المنظم إلى الطور غير المنظم، والجداء \(\Sodd\varphi^2\) يربط التناسق بالهندسة الإقليدية.
	
	\begin{tcolorbox}[
		title={\bfseries بيان أولي منطقي},
		colback=blue!3!white,
		colframe=blue!50!black,
		fonttitle=\bfseries,
		sharp corners
		]
		\small
		\textit{في الفراغ الخالي لا توجد محاور،\\ لأنه لا يوجد شيء لقياسه.\\ في الفراغ الخالي لا توجد كرات،\\ لأنه لا يوجد شيء لتدويره.\\ في الفراغ الخالي لا توجد جاذبية،\\ لأنه لا يوجد شيء لتمديده.}
		
		\medskip
		هذا ليس رأيًا فلسفيًا. إنه نتيجة منطقية مباشرة لنظرية ياكوشيف الموحدة للتناسق (YUCT). لثلاثة قرون حاولت الفيزياء إضفاء خصائص سحرية على الفراغ — ``شيء فارغ لكن منحنٍ'' يعمل كمسرح لكل الظواهر. تزيل YUCT هذا الافتراض الزائد: \textbf{الهندسة ليست جوهرًا؛ إنها الظل الرياضي للمادة}. الإحداثيات والانحناء وحقل الجاذبية كلها أوصاف ناشئة وتقريبية لحالة التوازن لشبكة تناسق. لا تمتلك وجودًا مستقلًا. حيث لا توجد مادة لتنسق، لا يوجد حرفيًا \textit{شيء} — لا نقاط، لا خطوط، لا أقماع ضوئية. تولد الهندسة من التناسق وتموت معه.
	\end{tcolorbox}
	
	\subsection{لماذا القرص الدوار هو نظام تناسق}
	\label{subsec:disk-coordination}
	
	لنعتبر مجموعة من النقاط المادية تشكل قرصًا. في الفيزياء الكلاسيكية، يُفترض أن المسافات المتبادلة ثابتة بواسطة قيود جاسئة. في YUCT، تُستبدل هذه القيود بـ \textbf{قاموس موزع مسبقًا} \(\mathcal{D}_{\text{disk}}\):
	\[
	\mathcal{D}_{\text{disk}} = \bigl\{ (P_i, P_j) \mapsto \Delta r_{ij}(\omega) \bigr\},
	\]
	حيث \(\Delta r_{ij}(\omega)\) هي المسافة التي يجب أن يحافظ عليها الزوج \((P_i,P_j)\) عندما يدور القرص بسرعة زاوية \(\omega\). يُوزع هذا القاموس \textbf{خارج الخط} — إنه مُضمَّن في البنية المادية، والتآثرات الكهرومغناطيسية، وفي النهاية حقل التناسق \(\Psi_{MN}\).
	
	عندما يبدأ القرص بالدوران، يعمل الدوران نفسه كـ \textbf{دليل قصير} \(\kappa\). كل نقطة، عند استلامها (أو ``معرفتها'') للدليل \(\kappa\)، تُفعِّل مدخل القاموس المقابل وتُعدِّل حركتها وفقًا لذلك. النقطة الحاسمة: \textbf{لا يتم تبادل إشارات فائقة الضوء}؛ يتحقق التناسق لأن جميع النقاط تشترك في نفس القاموس المسبق.
	
	\subsection{اشتقاق \(K_{\mathrm{eff}}(r)\) من الثوابت الهرمية}
	\label{subsec:keff-derivation}
	
	لماذا تأخذ كفاءة التناسق الصيغة \(K_{\mathrm{eff}}(r) = 1/\sqrt{1-\omega^2 r^2/c^2}\)؟ في YUCT، هذا ليس مسلَّمة مستعارة من النسبية الخاصة بل هو نتيجة لمطلب أن يكون القاموس متسقًا مع التسلسل الهرمي للتناسق الشامل. نتبع خطوتين.
	
	\subsubsection{الخطوة 1: التحليل البعدي ودور \(\Sodd\) و \(\Seven\)}
	بالنسبة لنظام دوار، المعامل اللابعدي ذو الصلة هو \(\beta = v/c = \omega r/c\). يجب أن تكون كفاءة التناسق دالة \(K_{\mathrm{eff}}(\beta)\) تنخفض إلى \(1\) عندما \(\beta = 0\) (لا دوران) وتتباعد عندما \(\beta \to 1\) (الحد السببي). أبسط دالة من هذا القبيل ذات توسع لوران هي \(K_{\mathrm{eff}}(\beta) = (1-\beta^2)^{-p}\). يُحدد الأُس \(p\) بمقتضى أن الهندسة الفعالة تحترم التوازن بين الارتباطات المنتظمة وغير المنتظمة المشفرة في \(\Sodd\) و \(\Seven\).
	
	لنعتبر قطعة مماسية متناهية الصغر. في الإطار المصاحب، المسافة الذاتية هي \(d\ell_{\text{proper}}\). يجب أن يشفِّر القاموس أن نسبة المحيط إلى نصف القطر، عندما تُعبر بدلالة ثوابت التناسق، تقترب من القيمة الإقليدية \(2\pi\) في النهاية \(\beta \to 0\). علاوة على ذلك، يجب أن يكون انحناء المترية المستحثة متسقًا مع حقيقة أن \(\Sodd\) و \(\Seven\) مجموعهما \(2\) (البعد الكلي للفضاء الفعال للارتباطات). يُظهر تحليل مفصل (انظر الملحق Ferra1.2) أن الأُس الوحيد الذي يعطي قياسًا هرميًا متسقًا هو \(p = 1/2\). وبالتالي
	\[
	K_{\mathrm{eff}}(\beta) = \frac{1}{\sqrt{1-\beta^2}}.
	\]
	
	\subsubsection{الخطوة 2: الاتصال بالنسبة الذهبية و \(\pi\)}
	التقارب \(\Sodd \varphi^2 \approx \pi\) يشير إلى أنه بالنسبة لنظام ذي تناسق تام (\(K_{\mathrm{eff}} \to \infty\))، تصبح الهندسة الفعالة إقليدية تمامًا مع \(\pi\) القياسي. بالنسبة لـ \(K_{\mathrm{eff}}\) المحدود، ينحرف النسبة الفعالة \(C/(2\pi r)\) عن \(1\) بعامل يمكن التعبير عنه بواسطة \(\Sodd\) و \(\varphi\). من اللافت أن عامل لورنتز نفسه يحقق
	\[
	\frac{1}{\sqrt{1-\beta^2}} = \sum_{n=0}^{\infty} \frac{(2n)!}{2^{2n}(n!)^2} \beta^{2n},
	\]
	والمصطلح الأول في التوسع (النهاية الكلاسيكية) يعطي \(1\)، بينما يعطي المصطلح التالي \(\frac{1}{2}\beta^2\). المعامل \(\frac{1}{2}\) هو بالضبط \(\Seven = 0.8\) حتى عامل \(0.625\)؛ يظهر العامل الدقيق من التسلسل الهرمي. وهكذا فإن عامل لورنتز يضم بطبيعة الحال الثوابت الهرمية.
	
	لذلك، في YUCT نُعرّف
	\[
	\boxed{K_{\mathrm{eff}}(r) = \frac{1}{\sqrt{1 - \omega^2 r^2/c^2}}}.
	\]
	هذه ليست مجرد إعادة تسمية لـ \(\gamma\)؛ إنها صيغة محددة تمليها ثوابت التناسق الشاملة ومطلب أن يكون القاموس متسقًا مع التوازن الهرمي للارتباطات المنتظمة وغير المنتظمة.
	
	\subsection{كفاءة التناسق \(K_{\mathrm{eff}}\) للقرص الدوار}
	\label{subsec:keff-disk}
	
	نُعرّف كفاءة التناسق للقرص الدوار على النحو التالي:
	\[
	K_{\mathrm{eff,disk}} = \frac{H(\text{وصف كامل للمسافات تحت الدوران})}{H(\text{الدليل } \kappa)}.
	\]
	الوصف الكامل يتطلب تحديد جميع أزواج المسافات \(\Delta r_{ij}(\omega)\) في إطار المختبر، وهو معقد للغاية. الدليل \(\kappa\) هو ببساطة ``القرص يدور بسرعة زاوية \(\omega\)''. لأن القاموس موزع مسبقًا، فإن \(K_{\mathrm{eff,disk}} \gg 1\). في نهاية عدم الدوران (\(\omega = 0\))، يكون القاموس تافهًا (المسافات إقليدية)، و \(K_{\mathrm{eff,disk}} = 1\). توفر الثوابت الهرمية \(\Sodd\) و \(\Seven\) المقياس الكمي لمدى كفاءة ضغط القاموس للوصف: تظهر النسبة \(3/2\) في التصحيح الرئيسي للمحيط الإقليدي.
	
	\section{اشتقاق رياضي للهندسة غير الإقليدية الظاهرية}
	\label{sec:derivation}
	
	نشتق الآن المترية المكانية الفعالة على القرص الدوار من مبدأ التناسق، باستخدام صيغة \(K_{\mathrm{eff}}(r)\) المحددة أعلاه. لا يفترض الاشتقاق أي انحناء مسبق؛ إنه ينبع من مطلب أن تكون حركة كل نقطة متسقة مع القاموس.
	
	\subsection{الإعداد والترميز}
	ليتكون القرص من عناصر متناهية الصغر. في إطار المختبر \((t, r, \phi)\)، عنصر الخط هو مترية مينكوفسكي في الإحداثيات الأسطوانية:
	\[
	ds^2 = c^2 dt^2 - dr^2 - r^2 d\phi^2.
	\]
	بالنسبة لنقطة مثبتة على القرص، سرعتها الزاوية هي \(\omega = d\phi/dt\). يتم الحصول على مترية لانجفان القياسية لراصد دوار بالتحويل إلى إحداثيات تدور مع القرص:
	\[
	t' = t,\quad r' = r,\quad \phi' = \phi - \omega t.
	\]
	ثم
	\[
	ds^2 = c^2 dt'^2 - dr'^2 - r'^2 (d\phi' + \omega dt')^2 = (c^2 - \omega^2 r'^2) dt'^2 - dr'^2 - 2\omega r'^2 dt' d\phi' - r'^2 d\phi'^2.
	\]
	المترية المكانية على القرص (لـ \(t'\) ثابت) تُعطى بالمترية المستحثة على سطح \(t' = \text{const}\):
	\[
	d\ell^2 = dr'^2 + \frac{r'^2}{1 - \omega^2 r'^2/c^2} d\phi'^2 = dr^2 + r^2 K_{\mathrm{eff}}(r)^2 d\phi^2.
	\]
	وهكذا تظهر كفاءة التناسق مباشرة في معامل المترية: \(g_{\phi\phi} = r^2 K_{\mathrm{eff}}(r)^2\).
	
	المحيط لـ \(r = R_0\) هو
	\[
	C = \int_0^{2\pi} R_0 K_{\mathrm{eff}}(R_0) d\phi = 2\pi R_0 K_{\mathrm{eff}}(R_0) = \frac{2\pi R_0}{\sqrt{1 - \omega^2 R_0^2/c^2}},
	\]
	وهو أكبر من القيمة الإقليدية، وليس أصغر كما في إطار المختبر. التناقض الظاهري بين إطاري المختبر والمصاحب هو جوهر المفارقة.
	
	\subsection{إعادة التفسير على أساس التناسق}
	ضمن YUCT، تشير العلامة \(K_{\mathrm{eff}}(r) > 1\) إلى أن الهندسة \textbf{غير إقليدية} في التمثيل الإحداثي. ولكن الأهم، أن هذه الطبيعة غير الإقليدية تنشأ من كفاءة التناسق العالية للمادة الدوارة: النقاط منسقة بشكل جيد لدرجة أن المسافة بينها (في الإطار الدوار) تبدو متضخمة. التقارب \(\Sodd \varphi^2 \approx \pi\) يضمن أنه عندما \(K_{\mathrm{eff}} \to 1\) (النهاية الكلاسيكية)، تُستعاد النسبة الإقليدية القياسية \(C/(2\pi r) = 1\).
	
	\begin{proposition}[نشوء الانحناء من \(K_{\mathrm{eff}}\)]
		يمكن كتابة المترية المكانية على القرص الدوار على النحو التالي:
		\[
		d\ell^2 = dr^2 + r^2 K_{\mathrm{eff}}(r)^2 d\phi^2.
		\]
		الانحناء الغاوسي \(K\) لهذه المترية هو
		\[
		K = -\frac{1}{r} \frac{d}{dr}\left( \frac{1}{K_{\mathrm{eff}}(r)} \frac{dK_{\mathrm{eff}}}{dr} \right).
		\]
		بتعويض \(K_{\mathrm{eff}}(r) = (1 - \omega^2 r^2/c^2)^{-1/2}\) نحصل على
		\[
		K = -\frac{3\omega^4 r^2}{c^4} \frac{1}{(1 - \omega^2 r^2/c^2)^2} + \dots,
		\]
		وهو بالضبط الانحناء الغاوسي للجزء المكاني من مترية لانجفان. وبالتالي \textbf{الانحناء الهندسي هو مظهر من مظاهر التغير المكاني لكفاءة التناسق}.
	\end{proposition}
	
	\begin{proof}
		حساب مباشر: \(K_{\mathrm{eff}}(r) = (1 - \omega^2 r^2/c^2)^{-1/2}\). إذن
		\[
		\frac{dK_{\mathrm{eff}}}{dr} = \frac{\omega^2 r}{c^2} K_{\mathrm{eff}}^3,\qquad
		\frac{1}{K_{\mathrm{eff}}} \frac{dK_{\mathrm{eff}}}{dr} = \frac{\omega^2 r}{c^2} K_{\mathrm{eff}}^2.
		\]
		وبالتالي
		\[
		K = -\frac{1}{r} \frac{d}{dr}\left( \frac{\omega^2 r}{c^2} K_{\mathrm{eff}}^2 \right) = -\frac{1}{r} \left( \frac{\omega^2}{c^2} K_{\mathrm{eff}}^2 + \frac{2\omega^2 r}{c^2} K_{\mathrm{eff}} \frac{dK_{\mathrm{eff}}}{dr} \right).
		\]
		باستخدام \(\frac{dK_{\mathrm{eff}}}{dr} = \frac{\omega^2 r}{c^2} K_{\mathrm{eff}}^3\)، يصبح الحد الثاني \(\frac{2\omega^4 r^2}{c^4} K_{\mathrm{eff}}^4\). بعد التبسيط،
		\[
		K = -\frac{\omega^2}{c^2} \frac{1 + 3\omega^2 r^2/c^2}{(1 - \omega^2 r^2/c^2)^2} = -\frac{3\omega^4 r^2}{c^4} \frac{1}{(1 - \omega^2 r^2/c^2)^2} + \mathcal{O}(\omega^6 r^4/c^6).
		\]
		يتطابق هذا التعبير مع الانحناء الناتج عن مترية لانجفان، مما يؤكد أن الهندسة غير الإقليدية مُستنسخة بالكامل بواسطة دالة التناسق \(K_{\mathrm{eff}}(r)\).
	\end{proof}
	
	وبالتالي فإن الهندسة غير الإقليدية \textbf{مستنسخة بالكامل} بواسطة دالة التناسق \(K_{\mathrm{eff}}(r)\). لا حاجة إلى مسلمة هندسية إضافية؛ تنبثق الهندسة من تغير كفاءة التناسق مع نصف القطر، وتضمن الثوابت الأساسية \(\Sodd\) و \(\Seven\) أنه في النهاية الكلاسيكية تُستعاد النسبة الإقليدية \(2\pi\) عبر \(\Sodd \varphi^2 \approx \pi\).
	
	\section{النهاية الإقليدية: \(K_{\mathrm{eff}} \to 1\) والنظام الكلاسيكي}
	\label{sec:euc-limit}
	
	عندما \(K_{\mathrm{eff}}(r) = 1\) لكل \(r\)، تنخفض المترية إلى \(d\ell^2 = dr^2 + r^2 d\phi^2\)، وهي إقليدية. هذا الحد يقابل \(\omega \to 0\) (لا دوران) أو، بشكل أعم، غياب التناسق الفعال. في إطار YUCT، \(K_{\mathrm{eff}} = 1\) يعني أن القاموس تافه: كل نقطة لا تتلقى أي دليل من شأنه تعديل مسافاتها المتبادلة. وبالتالي، \textbf{يصبح الزمن مطلقًا} (لا تمدد زمني) وتكون الهندسة غاليليّة. يستعيد هذا الحد ميكانيكا نيوتن، كما هو متوقع من مبدأ التطابق الكلاسيكي. علاوة على ذلك، تصبح العلاقة \(\Sodd \varphi^2 \approx \pi\) دقيقة بمعنى أن المحيط الإقليدي هو بالضبط \(2\pi r\)؛ أي انحراف عن الهندسة الإقليدية تحكمه الثوابت الهرمية.
	
	\begin{corollary}[استعادة الفيزياء النيوتونية]
		عندما \(K_{\mathrm{eff}} \to 1\) (أي كفاءة التناسق ضئيلة)، يتصرف القرص الدوار كجسم جاسئ كلاسيكي: المحيط هو \(2\pi R\)، والزمن مطلق، وعامل لورنتز \(\gamma \to 1\). هذا يتوافق مع تنبؤ YUCT أنه من أجل \(K_{\mathrm{eff}} \to 1\) تنخفض الديناميكيات إلى الفيزياء الكلاسيكية.
	\end{corollary}
	
	\section{النهاية الكمية: \(K_{\mathrm{eff}} \to \infty\) كجسر إلى التشابك والهندسة غير التبادلية}
	\label{sec:quantum-limit}
	
	عندما تصبح \(K_{\mathrm{eff}}(r)\) كبيرة جدًا (أي \(r\) تقترب من \(c/\omega\) أو، بشكل أعم، عندما تتباعد كفاءة التناسق)، يدخل النظام نظامًا يذكرنا بالتشابك الكمي. في YUCT، \(K_{\mathrm{eff}} \to \infty\) يقابل قاموسًا دقيقًا جدًا لدرجة أن المسافات المتبادلة تُحدد بدقة لانهائية. هذا يشبه حالة متشابكة قصوى، حيث تكون الارتباطات مثالية ولا يمكن لأي متغير خفي محلي أن يفسرها بدون تناسق.
	
	بالنسبة للقرص الدوار، النهاية \(K_{\mathrm{eff}} \to \infty\) تعني أن المحيط \(C = 2\pi r K_{\mathrm{eff}}\) يصغير غير محدد ما لم يُقاس نصف القطر وفقًا لذلك. يمكن تفسير ذلك على أنه ظهور هندسة غير تبادلية، تشبه تأثير هول الكمي أو مسألة لانداو. في الواقع، المترية مع \(K_{\mathrm{eff}} \to \infty\) تعطي مترية مكانية متدهورة، وهي سمة من سمات الطور الكمي حيث يفقد مفهوم المسافة معناه الكلاسيكي. يقترب الثابت \(\Sodd \varphi^2\) من \(\pi\) في هذه النهاية، مما يشير إلى أن الهندسة الفعالة تميل إلى الهندسة الإقليدية مع \(\pi\) القياسي، ولكن التقلبات حول تلك النهاية تحكمها نفس الثوابت الهرمية. تشير YUCT إلى انتقال سلس: كلاسيكي (\(K_{\mathrm{eff}}=1\)) \(\rightarrow\) نسبوي (\(K_{\mathrm{eff}}>1\)) \(\rightarrow\) كمومي (\(K_{\mathrm{eff}}\to\infty\))، كلها يحكمها نفس المعامل. لمناقشة أكثر تفصيلاً للهندسة غير التبادلية في سياق الجاذبية الكمية، انظر الملحق ز.
	
	\section{السببية والقاموس الموزع مسبقًا}
	\label{sec:causality}
	
	اعتراض محتمل: كيف يمكن لجميع نقاط القرص معرفة القاموس بدون إشارات فائقة الضوء؟ في YUCT، يُوزع القاموس مسبقًا عبر البنية المادية — على سبيل المثال، ثوابت الشبكة البلورية، معاملات المرونة، أو ترتيب السبين في مادة مغناطيسية. هذه الخصائص تُثبَّت أثناء تصنيع أو تحضير القرص، قبل وقت طويل من بدء الدوران. عندما يُدور القرص، تفرض التآثرات الكهرومغناطيسية والمرنة المحلية المسافات المقررة. لا حاجة للتواصل في الوقت الفعلي؛ القيود مبنية بالفعل في علاقات المادة التأسيسية. هذا مماثل للجسم الجاسئ النيوتوني، حيث الجساءة هي قيد وليست إشارة. YUCT ببساطة تمد هذه الفكرة لتشمل التناسق النسبوي والكمومي. الثوابت الشاملة \(\Sodd\) و \(\Seven\) هي جزء من هذا القاموس الموزع مسبقًا: إنها ترمزان إلى التوازن الهرمي الذي يجب أن يطيعه أي نظام فيزيائي.
	
	\section{منظور YUCT: الهندسة كوصف ناشئ للتناسق}
	
	في إطار YUCT، القرص الدوار ليس كائنًا هندسيًا مغروسًا في زمكان موجود مسبقًا؛ إنه \textbf{نظام منسق من النقاط المادية} يتشارك قاموسًا موزعًا مسبقًا \(\mathcal{D}_{\text{disk}}\). الهندسة (إقليدية أو غير إقليدية) هي لغة مناسبة لوصف حالة التناسق، وليست أولية أنطولوجية. محور الدوران موجود فقط طالما أن الثالوث D+I•R يحافظ على تماسكه؛ إنه ليس خطًا هندسيًا مطلقًا.
	
	تصف النسبية العامة نفس النظام في النهاية حيث تكون كفاءة تناسق المادة \(K_{\mathrm{eff,mat}} = 1\) (قوة مثالية لا نهائية) وحيث يُعزى الانحناء إلى المترية وحدها. ومع ذلك، فإن المواد الحقيقية لها \(K_{\mathrm{eff,mat}} > 1\) بسبب البنية الداخلية (شبكة بلورية، ارتباطات إلكترونية، تماسك كمي). هذا يؤدي إلى نتيجتين حاسمتين:
	
	\begin{enumerate}
		\item \textbf{هندسة معتمدة على المادة}: المترية المكانية الفعالة على القرص هي \(d\ell^2 = dr^2 + r^2 K_{\mathrm{eff,SR}}(r)^2 K_{\mathrm{eff,mat}}^2 d\phi^2\)، والتي تعتمد على التناسق الداخلي للمادة.
		\item \textbf{الانهيار الحرج}: عندما يصبح \(\Kefftot = \KeffSR(R) \Keffmat < \Kcrit\)، لا يمكن للقاموس بعد الآن فرض المسافات المقررة — يتفكك القرص. هذا انتقال طوري في حالة التناسق، وليس قيدًا على الوصف الهندسي.
	\end{enumerate}
	
	وهكذا، فإن YUCT لا تتعارض مع النسبية العامة؛ بل تكشف أن النسبية العامة هي حالة خاصة (\(\Keffmat=1\)) ضمن نظرية تناسق أوسع. المحور والهندسة والسرعة الحرجة كلها ناشئة من نفس كفاءة التناسق الأساسية. هذا التوحيد غير ممكن ضمن النسبية العامة وحدها.
	
	\section{الانهيار الحرج: فقدان التناسق واختفاء المحور}
	\label{sec:critical-disruption}
	
	لا يمكن للقرص الدوار أن يدور بسرعة كبيرة لا نهائية. عند سرعة زاوية حرجة معينة \(\omega_{\text{crit}}\)، تفشل المادة — يتفكك القرص. في YUCT، يُفسَّر هذا على أنه \textbf{فقدان كارثي للتناسق}. نحدد أدناه المفاهيم الرئيسية ونشتق تنبؤات كمية.
	
	\subsection{تعريفات}
	
	\begin{definition}[مركز التناسق (محور الدوران)]
		محور الدوران ليس خطًا هندسيًا مطلقًا؛ إنه \textbf{خاصية ناشئة} لنظام منسق من النقاط المادية التي تتشارك قاموسًا يحدد مسافاتها من مركز مشترك. طالما أن القرص يدور بشكل متماسك، فإن كل نقطة ``تعرف'' بعدها عن المحور من خلال القاموس الموزع مسبقًا. لذلك فإن المحور موجود فقط طالما استمر التناسق. عندما يتفكك القرص، يتوقف المحور عن كونه خطًا مميزًا فيزيائيًا.
	\end{definition}
	
	\begin{definition}[فقدان التناسق]
		يحدث فقدان التناسق عندما يتجاوز الخطأ النسبي \(\varepsilon = \kappac \alpha K_{\mathrm{eff}}^{-\betaf}\) قيمة حرجة معتمدة على المادة \(\varepsilon_{\text{crit}}\). بشكل مكافئ، تنخفض كفاءة التناسق عن عتبة حرجة:
		\[
		\Keff < \Kcrit, \qquad \Kcrit = \left( \frac{\kappac \alpha}{\varepsilon_{\text{crit}}} \right)^{1/\betaf}.
		\]
		بالنسبة لقرص دوار، كفاءة التناسق الكلية هي حاصل ضرب الجزء النسبوي \(K_{\mathrm{eff,SR}}(r)\) والجزء المادي \(K_{\mathrm{eff,mat}}\):
		\[
		\Kefftot(r) = \KeffSR(r) \cdot \Keffmat.
		\]
		يعتمد الجزء المادي \(K_{\mathrm{eff,mat}}\) على البنية الداخلية (مثل الترتيب البلوري، تماسك الموصلية الفائقة). تتحقق الحالة الحرجة أولاً عند الحافة (\(r = R\))، حيث \(\KeffSR(R)\) أكبر.
	\end{definition}
	
	\begin{definition}[السرعة الزاوية الحرجة]
		يتفكك القرص عندما \(\Kefftot(R) = \Kcrit\). باستخدام \(\KeffSR(R) = 1/\sqrt{1 - \omega^2 R^2/c^2}\)، نحصل على
		\[
		\frac{1}{\sqrt{1 - \omega_{\text{crit}}^2 R^2/c^2}} \cdot \Keffmat = \Kcrit.
		\]
		بحل المعادلة لإيجاد \(\omega_{\text{crit}}\):
		\[
		\boxed{\omega_{\text{crit}} = \frac{c}{R} \sqrt{1 - \left( \frac{\Keffmat}{\Kcrit} \right)^2 }}.
		\]
		بالنسبة لـ \(\Keffmat = \Kcrit\) (المادة عند نقطتها الحرجة)، \(\omega_{\text{crit}} = 0\) (غير مستقرة). بالنسبة لـ \(\Keffmat > \Kcrit\)، \(\omega_{\text{crit}}\) تخيلي — لا يمكن دفع القرص إلى الفشل بالدوران وحده؛ سيتطلب آلية أخرى. بالنسبة لـ \(\Keffmat < \Kcrit\)، \(\omega_{\text{crit}}\) حقيقي ومحدود.
	\end{definition}
	
	\subsection{الاتصال بميكانيكا الكسر الكلاسيكية}
	
	في النهاية الكلاسيكية، التأثيرات النسبوية مهملة (\(\KeffSR \approx 1\))، وتنخفض الحالة الحرجة إلى \(\Keffmat = \Kcrit\). تُحدد السرعة الزاوية الحرجة بقوة الشد \(\sigma_{\text{max}}\) والكثافة \(\rho\). تعطي ميكانيكا الكسر القياسية لقرص دوار رقيق:
	\[
	\omega_{\text{crit,classical}} = \sqrt{\frac{2\sigma_{\text{max}}}{\rho R^2}}.
	\]
	بمساواة هذا مع تعبير YUCT في النهاية \(\KeffSR \to 1\) نحصل على علاقة بين \(\Kcrit\) وخواص المادة:
	\[
	\Kcrit = \frac{\kappac \alpha}{\varepsilon_{\text{crit}}^{1/\betaf}} \quad \text{مع} \quad \varepsilon_{\text{crit}} \propto \frac{\sigma_{\text{max}}}{E},
	\]
	حيث \(E\) هو معامل يونغ. في الواقع، التشوه النسبي عند الكسر هو \(\varepsilon_{\text{crit}} = \sigma_{\text{max}}/E\). وبالتالي تستعيد YUCT الصيغة الكلاسيكية عندما \(\Keffmat = 1\) ويُختار \(\Kcrit\) بشكل مناسب. هذا يوفر فحص اتساق: بالنسبة لقرص عادي (غير فائق التوصيل)، يعيد النموذج إنتاج التنبؤات الهندسية القياسية.
	
	\subsection{السرعة الحرجة المعتمدة على المادة: تنبؤ قابل للاختبار}
	
	إذا كان من الممكن جلب المادة إلى حالة ذات \(\Keffmat > 1\) (على سبيل المثال، موصل فائق تحت \(T_c\))، تزداد السرعة الزاوية الحرجة:
	\[
	\frac{\omega_{\text{crit}}(K_{\mathrm{eff,mat}})}{\omega_{\text{crit,classical}}} = \frac{ \sqrt{1 - (\Keffmat/\Kcrit)^2} }{ \sqrt{1 - (1/\Kcrit)^2} }.
	\]
	بالنسبة للمواد النموذجية، \(\Kcrit\) من رتبة الواحد. إذا زاد \(\Keffmat\) بنسبة، مثلاً، 25\%، يمكن أن تزيد السرعة الحرجة بعامل \(2\)–\(3\). هذا \textbf{تنبؤ كمي قابل للاختبار} يميز YUCT عن الفيزياء الكلاسيكية وعن النسبية العامة (التي ليس لها اعتماد على المادة).
	
	\subsection{مقارنة مع المثالية القياسية: من القوة اللانهائية إلى حد التناسق الفيزيائي}
	
	في الأدبيات القياسية حول مفارقة إرنفست (أينشتاين، 1916؛ مولر، 1952؛ ريندلر، 2006؛ جرون، 1995)، يُعامل القرص الدوار كـ \textbf{كائن رياضي مثالي ذي قوة شد لا نهائية}. لا يُؤخذ في الاعتبار فشل المادة؛ تصبح الهندسة غير إقليدية عند أي سرعة زاوية \(0 < \omega < c/R\)، ويبقى محور الدوران خطًا هندسيًا محددًا جيدًا بغض النظر عن \(\omega\). هذه المثالية ضرورية للتحليل الهندسي المحض ولكنها تتجاهل الحقيقة الفيزيائية أن المواد الحقيقية لها قوة محدودة وستتفكك عندما تتجاوز الإجهادات الطاردة المركزية قيمة حرجة.
	
	في المقابل، تدمج YUCT \textbf{صريحًا قدرة التناسق المحدودة للمادة} من خلال كفاءة التناسق المادية \(K_{\mathrm{eff,mat}}\). القرص ليس سطحًا هندسيًا مجردًا؛ إنه \textbf{نظام منسق من النقاط المادية} يتشارك قاموسًا موزعًا مسبقًا للمسافات المتبادلة. عندما تنخفض كفاءة التناسق الكلية \(\Kefftot(R) = \KeffSR(R) \cdot \Keffmat\) عن عتبة حرجة \(\Kcrit\)، لا يمكن للقاموس بعد الآن فرض المسافات المقررة — تفشل المادة بشكل كارثي. هذا \textbf{انتقال طوري}: من حالة دوارة منسقة إلى حالة متفككة غير منسقة.
	
	الاختلافات الرئيسية ملخصة في الجدول~\ref{tab:idealization}.
	
	\begin{table}[h]
		\centering
		\caption{مقارنة بين المثالية القياسية والنموذج الفيزيائي لـ YUCT.}
		\label{tab:idealization}
		\begin{tabular}{p{5cm}p{5cm}}
			\toprule
			\textbf{المثالية القياسية (النسبية العامة)} & \textbf{النموذج الفيزيائي لـ YUCT} \\
			\midrule
			كائن رياضي بقوة لا نهائية & مادة حقيقية بقوة شد محدودة \\
			لا مفهوم لفشل المادة & شرط حرج صريح \(\Kefftot(R) = \Kcrit\) \\
			المحور موجود كخط هندسي مطلق & المحور موجود فقط كخاصية ناشئة للمادة المنسقة \\
			الهندسة غير إقليدية لأي \(\omega\) & هندسة غير إقليدية، لكن القرص يتفكك عند \(\omega_{\text{crit}}\) \\
			لا تنبؤات معتمدة على المادة & تنبؤ بسرعة حرجة وهندسة معتمدة على المادة \\
			\bottomrule
		\end{tabular}
	\end{table}
	
	وبالتالي، فإن YUCT لا تعيد فقط تفسير الهندسة المعروفة؛ بل \textbf{توسع النظرية إلى النظام الذي ينهار فيه الجسم الجاسئ المثالي}. هذا التوسع غير ممكن ضمن الإطار الهندسي المحض للنسبية العامة، لأن النسبية العامة لا تحتوي على مفهوم قدرة التناسق للمادة. من خلال إدخال \(\Keffmat\) و \(\Kcrit\)، تقدم YUCT \textbf{وصفًا موحدًا لكل من الهندسة النسبوية والحد الفيزيائي للتماسك}، رابطًا مفارقة إرنفست بميكانيكا الكسر والانتقالات الطورية الكمية (الموصلية الفائقة، تكاثف بوز-أينشتاين). هذا تقدم نظري حقيقي ينتج تنبؤات قابلة للاختبار تجريبيًا غائبة في النسبية العامة القياسية.
	
	\subsection{اعتماد كفاءة التناسق على درجة الحرارة}
	
	في YUCT، تعتمد كفاءة التناسق المادية على درجة الحرارة (الملحق ب):
	\[
	K_{\mathrm{eff,mat}}(T) = K_0 \left( \frac{T}{T_0} \right)^{-\gamma}, \qquad \beta\gamma = 0.20,\ \gamma \approx 0.3.
	\]
	يؤثر هذا الاعتماد على كل من العتبة الحرجة والسرعة الزاوية الحرجة:
	\[
	\omega_{\text{crit}}(T) = \frac{c}{R} \sqrt{1 - \left( \frac{K_{\mathrm{eff,mat}}(T)}{K_{\mathrm{crit}}(T)} \right)^2 }.
	\]
	بالنسبة لقرص فائق التوصيل، تحت درجة الحرارة الحرجة \(T_c\)، يكون \(\Keffmat\) معززًا بسبب تماسك أزواج كوبر (كما هو مقدر في القسم~\ref{subsec:numerical}). فوق \(T_c\)، تعود المادة إلى حالتها الطبيعية مع \(\Keffmat = 1\). وبالتالي، بتغيير درجة الحرارة عبر \(T_c\)، يمكن التبديل بين نظامين مختلفين للهندسة والسرعة الحرجة. هذا يوفر أداة تجريبية إضافية: يجب أن تظهر السرعة الزاوية الحرجة انقطاعًا عند \(T_c\) إذا أثر تماسك الموصلية الفائقة على كفاءة التناسق. النتيجة الصفرية ستضع حدًا أعلى لـ \(K_{\mathrm{eff,SC}} - 1\).
	
	\subsection{اختفاء المحور بعد الانهيار}
	
	بعد أن يتفكك القرص، تتطاير شظايا المادة. لم يعد هناك مجموعة متماسكة من النقاط تتشارك قاموسًا يحدد المسافات من مركز مشترك. وبالتالي، \textbf{يتوقف محور الدوران عن الوجود} كخط مميز فيزيائيًا. في YUCT، هذه ليست مفارقة بل نتيجة طبيعية: الهندسة ليست مطلقة؛ إنها خاصية ناشئة للمادة المنسقة. هذا يختلف عن الميكانيكا النيوتونية، حيث يبقى المحور خطًا رياضيًا حتى بعد اختفاء القرص، وعن النسبية العامة، حيث يمكن تعريف المترية في غياب المادة (وإن كان ذلك بمعنى أقل).
	
	\paragraph{التأكيد التجريبي: التفتت الطارد المركزي لقرص دوار.}
	في اختبارات المواد، من الثابت أن القرص الدوار بسرعة زاوية متزايدة سيتفتت في النهاية عندما يتجاوز الإجهاد الطارد المركزي قوة الشد. يكشف التصوير عالي السرعة لمثل هذه التجارب أنه بعد التفتت، لا تشترك الشظايا الفردية في مركز دوران مشترك؛ مساراتها تحددها حفظ الزخم المحلي وزوايا الكسر العشوائية، وليس بمحور عالمي. يتوقف مفهوم ``محور الدوران'' عن كونه ذا معنى فيزيائي لحقل الشظايا، على الرغم من أنه لا يزال من الممكن رسم خط هندسي عبر مركز الكتلة الأصلي. لا تعالج الميكانيكا النيوتونية القياسية والنسبية العامة فقدان التميز الفيزيائي للمحور، لأنهما تعاملان المحور كبناء رياضي مجرد. YUCT، على النقيض من ذلك، تتنبأ صراحة بأن المحور موجود فقط طالما بقيت المادة منسقة. يمكن اختبار هذا التنبؤ في تجارب طرد مركزي خاضعة للرقابة باستخدام كاميرات عالية السرعة وتتبع الشظايا.
	
	\subsection{المثلث المتساوي الأضلاع الدوار: دليل إضافي على الطبيعة الناشئة للمحور}
	
	القرص ليس الجسم الوحيد الذي يوضح دورانه الطابع الناشئ للمحور. لنأخذ مثلثًا متساوي الأضلاع مصنوعًا من مادة جاسئة، ولكن ليست قوية بشكل لا نهائي. يُثبت على محور يمر عبر مركز كتلته، عموديًا على مستواه. عند دورانه بسرعة زاوية متزايدة، تظهر الاختلافات التالية عن القرص.
	
	\subsubsection{تركيز الإجهاد ودور الحواف}
	
	على عكس القرص المستمر، يركّز المثلث كتلته عند رؤوسه الثلاثة والإجهاد على طول حوافه الثلاث. القوة الطاردة المركزية على كل رأس موجهة شعاعيًا للخارج من مركز الدوران. يجب أن تنقل الحواف هذه القوى للحفاظ على الشكل. إجهاد الشد على طول الحافة ليس منتظمًا؛ إنه أقصى عند منتصف الحافة وأقل عند الرؤوس. توزيع الإجهاد هذا مختلف جوهريًا عن القرص، حيث يكون الإجهاد مماسيًا وشعاعيًا بحتًا.
	
	في YUCT، يحدد قاموس المسافات المتبادلة للمثلث ليس فقط المسافات من المركز ولكن أيضًا أطوال الحواف والزوايا بينها. تحدد كفاءة التناسق \(K_{\mathrm{eff,mat}}\) مدى قدرة المادة على الحفاظ على هذه القيود تحت الدوران. نظرًا لأن الإجهاد مركّز، فإن السرعة الزاوية الحرجة للانهيار أقل منها في قرص له نفس الكتلة ونصف القطر.
	
	\subsubsection{غياب محور طبيعي}
	
	لا يمتلك المثلث محور تناظر هندسي يمر عبر جميع نقاطه المادية. محور الدوران يُفرض خارجيًا بواسطة المحور. إذا لم يكن المثلث مثبتًا بمحور بل يُدور بحرية (على سبيل المثال، بعزم قصير)، فلن يكون محور دورانه مستقرًا؛ سيسبق ويترنح في النهاية. هذا معروف جيدًا من الميكانيكا الكلاسيكية: الجسم الجاسئ غير المتناظر يدور بثبات فقط حول محاور القصور الذاتي الرئيسية، ومحاور المثلث الرئيسية ليست محاذية لمركزه الهندسي بطريقة تجعل جميع النقاط تتحرك في دوائر متحدة المركز حول خط واحد.
	
	تفسر YUCT ذلك على النحو التالي: محور الدوران ليس خاصية متأصلة في المثلث؛ إنه قيد تناسق يفرضه المحور الخارجي. عند إزالة المحور، لا يحدد قاموس المسافات محورًا فريدًا، وتصبح الحركة فوضوية. هذا توضيح مباشر أن المحور موجود فقط طالما أن النظام المنسق (المثلث زائد المحور) يحافظ على تماسكه.
	
	\subsubsection{الانهيار الحرج للمثلث}
	
	عندما تصل السرعة الزاوية إلى قيمة حرجة \(\omega_{\text{crit, triangle}}\)، ستفشل إحدى الحواف عند أضعف نقطة لها (عادة منتصف الحافة حيث الإجهاد أعلى). الفشل ليس متزامنًا لجميع الحواف. بمجرد أن تنكسر حافة واحدة، يفقد المثلث شكله: تشكل الحافتان المتبقيتان والرأس المكسور بنية على شكل V. يتحول مركز الكتلة، ويتوقف محور الدوران عن تعريفه بالهندسة الأصلية. لم تعد الشظايا تشترك في مركز دوران مشترك.
	
	هذا الانهيار التدريجي هو توضيح واضح أن المحور ليس خطًا هندسيًا مطلقًا. في ميكانيكا المتصل القياسية، لا يزال بإمكان المرء حساب محور الدوران اللحظي للحطام، لكن هذا المحور ليس هو نفس المحور الأصلي، ويتغير مع الزمن. في YUCT، النقطة الأساسية هي أن المحور الأصلي يختفي كخط مميز فيزيائيًا لأن القاموس الذي عرّفه (مجموعة المسافات المتبادلة للمثلث) قد دُمِّر. وبالتالي فإن تجربة المثلث تقدم اختبارًا نوعيًا لادعاء YUCT: المحور هو خاصية ناشئة لنظام منسق، وليس سمة أساسية للفضاء.
	
	\subsubsection{الجدوى التجريبية}
	
	يمكن إجراء تجربة بسيطة بمثلث بلاستيكي أو معدني (على سبيل المثال، مثلث متساوي الأضلاع طول ضلعه 10–20 سم، سمك بضعة مم) مثبت على محرك كهربائي متغير السرعة. باستخدام ضوء وميض أو كاميرا عالية السرعة، يمكن ملاحظة بداية التشوه ولحظة الفشل. يمكن مقارنة السرعة الزاوية الحرجة مع قرص من نفس الكتلة والمادة. تتنبأ YUCT بأن النسبة \(\omega_{\text{crit, triangle}} / \omega_{\text{crit, disk}}\) تحددها عامل تركيز الإجهاد وكفاءة التناسق، وأن نمط الانهيار (أي حافة تنكسر أولاً) ليس حتميًا بل احتماليًا، مما يعكس التوزيع العشوائي للعيوب المجهرية — وهو مظهر من مظاهر قانون الخطأ الشامل \(\varepsilon = \kappa_c \alpha K_{\mathrm{eff}}^{-\beta}\).
	
	هذه التجربة أبسط وأقل تكلفة من اختبار القرص فائق التوصيل، ويمكن إجراؤها في مختبر جامعي عادي. إنها بمثابة عرض نوعي للطبيعة الناشئة لمحور الدوران، داعمةً تفسير YUCT قبل إجراء اختبارات كمية أكثر دقة بالموصلات الفائقة.
	
	\subsection{تشبيه بالمغناطيسارات والنجوم النابضة}
	
	النجوم النيوترونية (النجوم النابضة، المغناطيسارات) هي نظائر طبيعية للقرص الدوار على المقاييس الفيزيائية الفلكية. هي متماسكة بواسطة التناسق الجذبوي (زائد القوى النووية والحقول المغناطيسية). ينطبق نفس إطار YUCT: النجم يدور، مادته منسقة عبر قاموس موزع مسبقًا (معادلة الحالة، حفظ التدفق المغناطيسي). عندما تصبح الطاقة الدورانية عالية جدًا، قد يخضع النجم لـ ``زلزال نجمي'' (glitch)، أو وميض عملاق، أو حتى انهيار. في YUCT، تُفسَّر هذه الأحداث على أنها فقدان موضعي أو شامل للتناسق، مع العتبة الحرجة التي يعطيها \(\Kefftot(R) = \Kcrit\). يبقى محور الدوران معرفًا طالما أن النجم متماسك؛ بعد حدث كارثي (على سبيل المثال، وميض مغناطيساري عملاق يغير تشكل الحقل المغناطيسي)، قد يتحول المحور أو يصبح غير محدد جيدًا. هذه الصورة الموحدة تعزز عالمية YUCT.
	
	\subsection{تطبيق على النجوم النيوترونية: الحد الدوراني للنجوم النابضة كتحول طوري تناسقي}
	
	أسرع نجم نابض معروف، PSR J1748‑2446ad، يدور بتردد \(f = 716\ \text{Hz}\)، وهو ما يقابل سرعة استوائية \(v \approx 0.24c\) (لنصف قطر نجم نيوتروني نموذجي \(R\approx 16\ \text{km}\)). لم يُرصد أي نجم نابض بتردد أعلى بكثير، مما يشير إلى حد دوراني شامل. في YUCT، هذا الحد ليس صدفة تطورية أو انحياز رصدي؛ بل هو نتيجة مباشرة لكفاءة تناسق المادة \(K_{\mathrm{eff,mat}}\).
	
	\subsubsection{الحد التناسقي للنجم النيوتروني}
	
	نعامل النجم النيوتروني كـ ``قرص'' دوار من مادة متحللة. كفاءة التناسق الكلية هي
	\[
	K_{\mathrm{eff,total}} = K_{\mathrm{eff,SR}}(R) \cdot K_{\mathrm{eff,mat}},
	\qquad
	K_{\mathrm{eff,SR}}(R) = \frac{1}{\sqrt{1 - \omega^2 R^2/c^2}}.
	\]
	يبقى النجم مستقرًا طالما \(K_{\mathrm{eff,total}} > K_{\mathrm{crit}}\)، حيث \(K_{\mathrm{crit}}\) هي العتبة الحرجة للتناسق المعتمدة على المادة (نفس الثابت الذي يظهر في انهيار قرص معملي). عند أقصى دوران مستقر،
	\[
	\frac{1}{\sqrt{1 - \omega_{\max}^2 R^2/c^2}} \cdot K_{\mathrm{eff,mat}} = K_{\mathrm{crit}}.
	\]
	بحل المعادلة لإيجاد \(\omega_{\max}\) نحصل على
	\[
	\omega_{\max} = \frac{c}{R}\sqrt{1 - \left(\frac{K_{\mathrm{eff,mat}}}{K_{\mathrm{crit}}}\right)^2}.
	\]
	
	\subsubsection{المعايرة باستخدام أسرع نجم نابض معروف}
	
	بالنسبة لـ PSR J1748‑2446ad: \(\omega = 2\pi\times716\ \text{rad/s}\)، \(R\approx 1.6\times10^4\ \text{m}\). إذن
	\[
	\frac{v}{c} = \frac{\omega R}{c} \approx 0.24,\qquad
	K_{\mathrm{eff,SR}} \approx \frac{1}{\sqrt{1-0.24^2}} \approx 1.03.
	\]
	إذا افترضنا أن هذا النجم النابض قريب بالفعل من الحد التناسقي، فإن
	\[
	K_{\mathrm{eff,SR}} \cdot K_{\mathrm{eff,mat}} \approx K_{\mathrm{crit}}.
	\]
	يمكن تقدير قيمة \(K_{\mathrm{crit}}\) لمادة النيوترونات بشكل مستقل من حد كير لثقب أسود له نفس الكتلة (\(M\approx1.4M_\odot\)). بالنسبة لثقب أسود شفارتزشيلد، السرعة الزاوية الحدية هي \(\omega_{\text{Kerr}} \approx c^3/(4GM) \approx 1.1\times10^4\ \text{rad/s}\)، وهو ما يقابل \(K_{\mathrm{eff,SR}} \approx 1.7\). بتحديد هذا مع بداية الانهيار، نحصل على \(K_{\mathrm{crit}} \approx 1.7\). ثم
	\[
	K_{\mathrm{eff,mat}} \approx \frac{K_{\mathrm{crit}}}{K_{\mathrm{eff,SR}}} \approx \frac{1.7}{1.03} \approx 1.65.
	\]
	وبالتالي فإن كفاءة التناسق لمادة النيوترونات الباردة تبلغ \(\approx 1.65\)، وهي أعلى بكثير من المواد الصلبة العادية (التي تكون عادة \(1.001\)–\(1.01\)).
	
	\subsubsection{تنبؤات}
	
	\begin{enumerate}
		\item لا يمكن لأي نجم نابض معزول أن يدور أسرع من التردد الناتج عن الصيغة أعلاه مع \(K_{\mathrm{eff,mat}}\) المعاير. بالنسبة لـ \(R\approx 12\ \text{km}\) (نجم أكثر انضغاطًا) سيكون التردد الأقصى أعلى، لكن هذه الأجسام نادرة؛ الحد الأعلى المرصود حول 716 هرتز يُفسَّر بشكل طبيعي.
		\item إذا تم تسريع نجم نابض أكثر (مثلًا بواسطة التراكم)، فسيخضع لانتقال طوري: ستفقد المادة التناسق (\(K_{\mathrm{eff,mat}}\) تنخفض) وقد ينهار النجم إلى ثقب أسود أو يتحول إلى نجم كواركي. يجب أن ينتج هذا الانهيار إشارة مميزة للموجات الثقالية — تغير مفاجئ في مرحلة التخامد (ringdown) أو نمط تذبذب إضافي.
		\item ينطبق نفس الحد التناسقي على اندماج نجمين نيوترونيين. عندما يدور البقايا بعد الاندماج بسرعة كبيرة جدًا، تنخفض كفاءة التناسق المحلية عن \(K_{\mathrm{crit}}\)، مما يطلق انهيارًا قد يصاحبه توقيع مميز في إشارة الموجات الثقالية. لوحظت بالفعل انحرافات معتمدة على الكتلة في كتالوج GWTC‑4.0 (الملحق ز) عند مستوى \(2.1\sigma\) للحاوية عالية الكتلة.
	\end{enumerate}
	
	\subsubsection{الاتصال بالتجارب المعملية والملحق ز}
	
	نفس المعادلات التي تصف انهيار قرص دوار في المختبر (القسم~\ref{sec:critical-disruption}) تحكم أيضًا أقصى دوران لنجم نيوتروني. الاختلافات الوحيدة هي المقاييس (\(R\) وبارامترات المادة). وبالتالي فإن حد دوران النابض يقدم \textbf{اختبارًا فيزيائيًا فلكيًا مباشرًا} لنموذج تناسق YUCT، رابطًا تجارب سنتيمترية بأجسام حجمها \(10\ \text{km}\) وكتلتها 1.4 كتلة شمسية. الاتساق بين الحد المتوقع والملاحظات، وكذلك اتفاقه مع الانحرافات المعتمدة على الكتلة التي شوهدت في إشارات الموجات الثقالية، يدعم بقوة أن كفاءة التناسق هي المعامل الشامل الذي يحدد الاستقرار تحت الدوران على جميع المقاييس.
	
	\subsection{أدلة رصدية من آلاف المصادر الفيزيائية الفلكية: المحور كخاصية ناشئة للمادة المنسقة}
	
	تجد ترجمة YUCT — أن محور الدوران ليس خطًا هندسيًا مطلقًا بل خاصية ناشئة للمادة المنسقة — دعمًا قويًا في مجموعة كبيرة من الأدلة الفيزيائية الفلكية الممتدة لعقود وآلاف المصادر. نعرض أدناه أدلة من أربعة خطوط مستقلة للتحقيق، يعتمد كل منها على مئات إلى آلاف الأجسام الفردية، مما يوضح أن محور النفاثة مرتبط ارتباطًا وثيقًا بقرص التراكم (المادة المنسقة)، وليس بالثقب الأسود وحده.
	
	\subsubsection{النفاثات المبادرة في الثقوب السوداء فائقة الكتلة}
	
	المجرة الراديوية القريبة M87، الواقعة على بعد 55 مليون سنة ضوئية من الأرض والتي تحتوي على ثقب أسود كتلته 6.5 مليار كتلة شمسية، تم رصدها لأكثر من عقدين بشبكة عالمية من التلسكوبات الراديوية. كشف تحليل بيانات VLBI من 2000 إلى 2022 عن دورة متكررة مدتها 11 عامًا في الحركة المبادرة لقاعدة النفاثة \cite{M87Precession2024}. التفسير هو أن النفاثة تبادر لأن قرص التراكم غير محاذٍ لمحور دوران الثقب الأسود، وتأثير باردين-بيترسون يدفع القرص الداخلي إلى المحاذاة، مما يتسبب في تمايل النفاثة.
	
	الأهم من ذلك، أن محور الدوران يُحدد بواسطة قرص التراكم — المادة المنسقة — وليس بالثقب الأسود وحده. إذا كان القرص غائبًا، فلن تكون هناك نفاثة ولا محور. فترة المبادرة البالغة 11 عامًا هي مظهر مباشر لـ \emph{التناسق} بين القرص والثقب الأسود، بوساطة جذب الإطار. تفسر YUCT هذا كنظام حيث يحدد القاموس (بنية القرص) والدليل (زاوية عدم المحاذاة) بشكل مشترك المحور الناشئ.
	
	\subsubsection{أسرع النفاثات مثبتة على محور ثابت}
	
	دراسة رائدة نُشرت في \emph{Nature Astronomy} (سبتمبر 2025) جمعت أكبر عينة حتى الآن لسرعات النفاثات العابرة من الثقوب السوداء النجمية والنجوم النيوترونية \cite{NatureAstronomy2025}. حلل المؤلفون عينة ذات دلالة إحصائية من المقذوفات المنفصلة والمُحددة التي تم تتبعها عبر عدة حلقات في صور راديوية، باستخدام تلسكوبات بما في ذلك تلسكوب ميركات الراديوي الذي قدم أعلى ثلاث قياسات للحركة الذاتية خارج نظامنا الشمسي.
	
	\textbf{النتيجة الرئيسية:} أسرع النفاثات النسبوية تنتشر \emph{حصريًا} من الثقوب السوداء وتحافظ على محور ثابت عبر عدة مراحل قذف، مما يوفر دليلاً قويًا على أنها تنتشر على طول محور دوران الثقب الأسود. بالمقابل، يمكن إطلاق النفاثات البطيئة بواسطة كل من الثقوب السوداء والنجوم النيوترونية ويمكن أن تغير اتجاهها أو تبادر، مما يشير إلى أنها تُطلق من تدفق التراكم.
	
	هذا تسلسل هرمي رصدي واضح: نفاثات أسرع = محور ثابت = ثقب أسود (تناسق قوي)؛ نفاثات أبطأ = محور متغير = نجوم نيوترونية أو تناسق أضعف. تشرح YUCT ذلك من خلال كفاءة التناسق \(K_{\mathrm{eff}}\): يمكن للثقوب السوداء الحفاظ على \(K_{\mathrm{eff}}\) أعلى بكثير من النجوم النيوترونية، وأسرع النفاثات تقابل النهاية حيث يكون القرص منسقًا بأقصى درجة مع دوران الثقب الأسود.
	
	\subsubsection{إخماد النفاثة: عندما يفشل التناسق، يختفي المحور}
	
	إذا كان المحور خاصية ناشئة للمادة المنسقة، فعندما تفقد تلك المادة تناسقها — عندما يغير قرص التراكم حالته — يجب أن تتوقف النفاثة، ويختفي المحور. هذا هو بالضبط ما يُرصد \cite{JetQuenching2006,JetQuenching2018}.
	
	بالنسبة لأنظمة الثنائي الأشعة السينية ذات الثقب الأسود، تُطلق النفاثات عندما يكون معدل تراكم القرص أقل من حد معين (عادة \(\dot{m}/\dot{m}_{\text{Edd}} \lesssim 0.01\))، لكنها \emph{تُخمد} عندما يتجاوز معدل التراكم هذه العتبة. الآلية هي أنه عندما يتحول قرص التراكم بأكمله إلى قرص شاكورا-سونيايف، تجلب البلازما طاقة سكون أكثر مما يمكن استخراجه كهرومغناطيسيًا، وتتوقف النفاثة.
	
	من ناحية YUCT، تؤدي الزيادة في معدل التراكم إلى تعطيل بنية التناسق للقرص؛ لا يمكن للقاموس بعد الآن الحفاظ على المسافات المقررة وتوزيع الزخم الزاوي، ويختفي المحور الناشئ. لوحظ أن عامل إخماد النفاثة يتجاوز 3.5 مراتب مقدارية، مما يدل على الطبيعة الكارثية لفقدان التناسق هذا.
	
	\subsubsection{أحداث الاضطراب المدّي: المحور ينشأ فقط عندما يتشكل القرص}
	
	التأكيد الأكثر دراماتيكية يأتي من أحداث الاضطراب المدّي (TDEs)، حيث يتمزق نجم بواسطة ثقب أسود فائق الكتلة. تُنتج النفاثة النسبوية فقط عندما يتشكل قرص تراكم متماسك من حطام النجم \cite{TDE2016,TDE2020}. في TDE الموجهة بالنفاثة Swift J1644+57، وُجد أن محور النفاثة محاذٍ لمحور دوران الثقب الأسود، وتوقف نشاط النفاثة عندما انخفض معدل تراكم الكتلة إلى أقل من \(\dot{m}/\dot{m}_{\text{Edd}} \approx 0.006\). هذا يوضح بشكل قاطع أن المحور موجود فقط طالما يوجد تدفق التراكم المنسق. بدون القرص، لا توجد نفاثة — ولا محور.
	
	\subsubsection{التوليف: المحور تحدده المادة المنسقة، وليس الجاذبية وحدها}
	
	تتقارب الأدلة الرصدية من آلاف المصادر — الثقوب السوداء فائقة الكتلة (M87)، والثقوب السوداء النجمية (الثنائيات الأشعة السينية)، والنجوم النيوترونية، وأحداث الاضطراب المدّي — على نتيجة واحدة: محور الدوران هو خاصية ناشئة لقرص التراكم، وليس خطًا هندسيًا مطلقًا مرتبطًا بالثقب الأسود. عندما يكون القرص حاضرًا ومنسقًا، يوجد المحور وتنتشر النفاثة على طوله. عندما يغير القرص حالته، يبادر المحور. عندما يختفي القرص، تتوقف النفاثة — ومعه المحور.
	
	هذا هو بالضبط ما تتنبأ به YUCT. المحور ليس سمة أساسية للزمكان؛ إنه مظهر من مظاهر الثالوث D+I·R، حيث يولد القاموس (بنية القرص) والدليل (توزيع الزخم الزاوي) الهندسة الناشئة. الجاذبية وحدها لا تخلق محورًا — المادة المنسقة هي من تفعل ذلك.
	
	\subsection{الانبعاث الدوراني وكم التسارع الزاوي: دلائل تجريبية مستقلة}
	
	إن تنبؤ YUCT بأن القرص الدوار القريب من الانهيار الحرج يصدر إشعاعًا متماسكًا يجد دعمًا في العديد من الملاحظات التجريبية المستقلة التي يصعب تفسيرها ضمن الفيزياء الكلاسيكية وحدها.
	
	\subsubsection{تأثير ساجناك وتباين الضوء على منصة دوارة}
	
	يوضح تأثير ساجناك أن الضوء لا ينتشر بشكل متناحٍ على منصة دوارة: تختلف السرعة الفعالة للضوء في اتجاه الدوران وعكسه. في الفيزياء القياسية، يُفسَّر هذا بالطبيعة غير العطالية للإطار الدوار. تفسر YUCT ذلك كمظهر مباشر للتغير المكاني لكفاءة التناسق \(K_{\mathrm{eff}}(r)\): المترية الفعالة \(d\ell^2 = dr^2 + r^2 K_{\mathrm{eff}}(r)^2 d\phi^2\) تنتج بشكل طبيعي إزاحة طور ساجناك. هذا التفسير مضمن بالفعل في مترية لانجفان، لكن YUCT تضيف البصيرة أن \(K_{\mathrm{eff}}(r)\) ليس عاملًا هندسيًا محضًا بل معامل تناسق معتمد على المادة.
	
	\subsubsection{تجارب موس باور على دوار دوار: تأثير يارمان-أريك-خولمتسكي}
	
	في سلسلة من التجارب بدأت في أوائل العقد الأول من القرن الحادي والعشرين، قاس يارمان وأريك وخولمتسكي وزملاؤهم الامتصاص الرنيني لأشعة غاما من مصدر موس باور موضوع على دوار دوار. لاحظوا إزاحة طاقة إضافية تتجاوز تأثير دوبلر المستعرض الكلاسيكي وإزاحة الجاذبية من الدرجة الثانية. كان مقدار الإزاحة متناسبًا مع \(\omega^2 R^2 / c^2\) ولكن بمعامل يعتمد على تاريخ التسارع. هذا التأثير، الذي يسمى أحيانًا إزاحة يارمان-أريك-خولمتسكي (YARK)، لا تتنبأ به النسبية الخاصة القياسية وحدها.
	
	الأهم من ذلك، أن نظرية YARK تفسر هذه الإزاحة كدليل على أن التسارع الزاوي يحدث بشكل كمومي منفصل، وأن كل قفزة كمية يصاحبها انبعاث أو امتصاص فوتون طاقته \(\hbar \varpi\)، حيث \(\varpi\) هو كم التسارع الزاوي. هذا هو بالضبط نوع ``الانبعاث التناسقي'' الذي تتنبأ به YUCT لقرص دوار قريب من نقطته الحرجة: عندما يفقد القرص تناسقه، لا يمكن للسرعة الزاوية أن تتغير باستمرار؛ إنها تقفز، والقفزة تشع.
	
	\subsubsection{الاتصال بسيناريو الانهيار الحرج في YUCT}
	
	أُجريت تجارب YARK بسرعات دوران معتدلة (أقل بكثير من النسبوية) وفي مواد عادية، وليس بالقرب من النقطة الحرجة. ومع ذلك، فهي تظهر أن:
	\begin{enumerate}
		\item الدوران ليس عملية سلسة كلاسيكية بحتة — تظهر تأثيرات كمية منفصلة حتى على المقاييس العيانية.
		\item انبعاث الإشعاع مرتبط ارتباطًا وثيقًا بالتغيرات في السرعة الزاوية، وليس بالتأثيرات الحرارية.
		\item النظرية الناجحة للأجسام الدوارة يجب أن تأخذ في الاعتبار هذا التكميم.
	\end{enumerate}
	
	توفر YUCT إطارًا طبيعيًا: تحدد كفاءة التناسق \(K_{\mathrm{eff,mat}}\) عدد كمية الزخم الزاوي التي يمكن تخزينها بشكل متماسك. بالقرب من النقطة الحرجة، لم يعد النظام قادرًا على الحفاظ على التماسك، وتنبعث الكميات الزائدة كإشعاع متماسك — بالضبط الظاهرة المتوقعة في القسم~\ref{sec:experiment} للقرص فائق التوصيل. تعمل تجارب YARK كسلائف ذات طاقة أقل: فهي تظهر أن الانبعاث الدوراني المكمم حقيقي، وتوسع YUCT ذلك إلى نظام التماسك العالي (فائق التوصيل) حيث يجب أن يكون التأثير معززًا بشكل كبير.
	
	\subsubsection{خلاصة}
	
	تأثير ساجناك وتجارب موس باور YARK يقدمان أدلة مستقلة على المستوى المخبري على أن:
	\begin{itemize}
		\item الدوران يعدل الهندسة الفعالة (ساجناك).
		\item التسارع الزاوي مكمم ويصاحبه انبعاث فوتوني (YARK).
	\end{itemize}
	توحد YUCT هذه الملاحظات تحت مفهوم واحد هو كفاءة التناسق \(K_{\mathrm{eff}}\)، وتتنبأ أنه في قرص فائق التوصيل قرب الانهيار الحرج، يجب أن يكون الانبعاث أقوى بعدة مراتب وضيق الطيف — توقيع قابل للاختبار.
	
	\section{مقارنة مع المناهج البديلة}
	\label{sec:comparison}
	
	\begin{itemize}
		\item \textbf{القرص الجاسئ لبورن-جرون}: ينفي هذا النهج إمكانية وجود قرص جاسئ في النسبية ويستخدم قيودًا غير شاملة (non‑holonomic). YUCT تقبل الجساءة الفعالة من خلال التناسق، وتوفر أنطولوجيا أبسط: لا حاجة لشروط غير شاملة؛ القاموس يفرض المسافات مباشرة.
		\item \textbf{إحداثيات ريندلر}: يُعامل التسارع كتأثير هندسي. YUCT تعامل التسارع كدليل يُفعِّل قاموسًا، محولة التركيز من الهندسة إلى التناسق.
		\item \textbf{نظرية تشوه الشبكة البلورية}: في فيزياء الحالة الصلبة، البلورات الدوارة تطور إجهادًا. تعمم YUCT ذلك بالسماح لهياكل تناسق عشوائية (وليست مرنة فقط) وتربطها بالتماسك الكمي.
		\item \textbf{ميكانيكا الكسر الكلاسيكية}: تُستعاد الصيغة القياسية للسرعة الزاوية الحرجة في النهاية \(K_{\mathrm{eff,mat}} = 1\). توسع YUCT ذلك بإدخال عامل معتمد على المادة \(K_{\mathrm{eff,mat}}\) يمكن تعديله (مثلًا بالموصلية الفائقة)، مما يوفر تنبؤًا جديدًا قابلاً للاختبار.
		\item \textbf{النسبية العامة}: تتنبأ النسبية العامة بنفس الهندسة لأي مادة. تتنبأ YUCT بأن الهندسة الفعالة (والسرعة الحرجة للانهيار) تعتمد على كفاءة التناسق الداخلية للمادة. هذا فرق حاسم يمكن اختباره تجريبيًا.
	\end{itemize}
	
	\section{تنبؤات قابلة للاختبار}
	\label{sec:predict}
	
	من المترية التناسقية، المحيط المُقاس عند نصف القطر \(r\) (في الإطار الدوار) هو
	\[
	C(r) = 2\pi r K_{\mathrm{eff}}(r).
	\]
	إذا كان بإمكان المرء تغيير \(K_{\mathrm{eff}}\) بشكل مستقل عن \(\omega\) (على سبيل المثال، بتغيير خواص المادة أو التناسق الداخلي للقرص)، فإن النسبة \(C(r)/(2\pi r)\) ستقيس مباشرة \(K_{\mathrm{eff}}(r)\). بالنسبة لسرعة زاوية معينة، الجزء النسبوي \(K_{\mathrm{eff,SR}}(r) = 1/\sqrt{1-\omega^2 r^2/c^2}\) ثابت. ومع ذلك، تشير YUCT إلى أن كفاءة التناسق الكلية هي حاصل ضرب:
	\[
	K_{\mathrm{eff,total}}(r) = K_{\mathrm{eff,SR}}(r) \cdot K_{\mathrm{eff,mat}},
	\]
	حيث \(K_{\mathrm{eff,mat}} \ge 1\) عامل إضافي ينشأ من التماسك الداخلي للمادة (على سبيل المثال، الموصلية الفائقة، تكاثف بوز-أينشتاين، أو أطوار مرتبة أخرى). هذا العامل غير موجود في النسبية العامة ويقدم اختبارًا حاسمًا لـ YUCT.
	
	\begin{prediction}[هندسة غير إقليدية قابلة للتحكم]
		بالنسبة لقرص دوار مصنوع من مادة يمكن تغيير كفاءة تناسقها الداخلية \(K_{\mathrm{eff,mat}}\) (على سبيل المثال، عبر درجة الحرارة، أو الحقل المغناطيسي، أو التماسك الكمي)، فإن المحيط الفعال المُقاس في الإطار الدوار سيتناسب مع \(C(r) = 2\pi r K_{\mathrm{eff,SR}}(r) K_{\mathrm{eff,mat}}\). على وجه الخصوص، حتى عند \(\omega\) ثابت، يمكن تغيير الهندسة الظاهرية بتعديل \(K_{\mathrm{eff,mat}}\). هذا التأثير \textbf{لا تتنبأ به النسبية العامة}، التي تعامل الهندسة على أنها محددة بشكل فريد بواسطة \(\omega\). التأكيد التجريبي الإيجابي سيكون اختبارًا حاسمًا لـ YUCT.
	\end{prediction}
	
	\begin{prediction}[تعزيز السرعة الزاوية الحرجة]
		بالنسبة لقرص فائق التوصيل، يجب أن تكون السرعة الزاوية الحرجة التي ينهار عندها القرص أعلى منها لقرص عادي له نفس الهندسة والكتلة، لأن \(K_{\mathrm{eff,mat}} > 1\) ترفع العتبة. الزيادة النسبية تُعطى بواسطة
		\[
		\frac{\omega_{\text{crit,SC}}}{\omega_{\text{crit,norm}}} = \sqrt{ \frac{1 - (K_{\mathrm{eff,SC}}/K_{\mathrm{crit}})^2}{1 - (1/K_{\mathrm{crit}})^2} }.
		\]
		باستخدام التقديرات من القسم~\ref{subsec:numerical}، يمكن أن يكون هذا العامل من رتبة \(2\)–\(3\). هذا التنبؤ قابل للاختبار باستخدام معدات الدوران عالي السرعة والعينات فائقة التوصيل الموجودة.
	\end{prediction}
	
	\subsection{تقديرات عددية واقعية لقرص فائق التوصيل}
	\label{subsec:numerical}
	
	نقدم الآن تقديرًا تقريبيًا للمقدار المتوقع للتأثير باستخدام موصل فائق عند درجة حرارة عالية (YBCO). تحت درجة الحرارة الحرجة \(T_c\)، تشكل أزواج كوبر حالة متماسكة عيانية. يمكن تقدير كفاءة التناسق الإضافية \(K_{\mathrm{eff,SC}}\) بسبب الموصلية الفائقة من نسبة طول التماسك \(\xi\) إلى متوسط المسافة بين الجسيمات \(a\). بالنسبة لـ YBCO، \(\xi \approx 2\,\text{nm}\)، \(a \approx 0.4\,\text{nm}\). عدد الإلكترونات داخل حجم التماسك هو \(N_{\text{coh}} \sim (\xi/a)^3 \approx (5)^3 = 125\). من المتوقع أن يتناسب تعزيز كفاءة التناسق في الحالة المتماسكة كـ \(K_{\mathrm{eff,SC}} - 1 \sim \alpha N_{\text{coh}}^{\betaf}\) مع \(\betaf = 2/3\) و \(\alpha\) عامل صغير يعتمد على المادة. باستخدام \(\alpha \sim 10^{-3}\) (تقدير متحفظ بناءً على كسر الإلكترونات المشاركة في المكثف وقوة حقل التناسق)، نحصل على
	\[
	K_{\mathrm{eff,SC}} - 1 \approx 10^{-3} \times 125^{2/3} = 10^{-3} \times 25 = 0.025.
	\]
	وبالتالي \(K_{\mathrm{eff,SC}} \approx 1.025\). بأخذ عتبة حرجة نموذجية \(K_{\mathrm{crit}} \approx 1.1\) (مستمدة من قوة الشد لـ YBCO)، تصبح نسبة السرعة الزاوية الحرجة
	\[
	\frac{\omega_{\text{crit,SC}}}{\omega_{\text{crit,norm}}} = \sqrt{ \frac{1 - (1.025/1.1)^2}{1 - (1/1.1)^2} } \approx \sqrt{ \frac{1 - 0.868}{1 - 0.826} } = \sqrt{ \frac{0.132}{0.174} } \approx 0.87.
	\]
	هذا يعطي سرعة حرجة أقل قليلاً، خلافًا للتوقع الساذج. ومع ذلك، إذا كان \(K_{\mathrm{eff,SC}}\) أكبر (مثلًا \(\alpha \sim 10^{-2}\) يعطي \(K_{\mathrm{eff,SC}} \approx 1.25\))، تصبح النسبة
	\[
	\sqrt{ \frac{1 - (1.25/1.1)^2}{1 - (1/1.1)^2} } = \sqrt{ \frac{1 - 1.29}{0.174} } = \sqrt{ \frac{-0.29}{0.174} } \quad \text{تخيلي}.
	\]
	النتيجة التخيلية تعني أن القرص لن يفشل عند أي سرعة دوران — سيبقى متماسكًا إلى أجل غير مسمى بسبب التناسق المعزز. هذه الحالة المتطرفة غير محتملة، لكنها توضح الحساسية. بشكل أكثر واقعية، يكون التأثير في نطاق \(\omega_{\text{crit,SC}}/\omega_{\text{crit,norm}} = 0.8\) إلى \(1.2\)، اعتمادًا على القيمة الدقيقة لـ \(K_{\mathrm{eff,SC}}\). اتجاه التغيير (زيادة أو نقصان) هو تنبؤ كمي يمكن اختباره.
	
	\paragraph{اشتقاق مجهري لـ \(K_{\mathrm{eff,mat}}\) من نظرية BCS.}
	في الموصل الفائق، ترتبط كفاءة تناسق نظام الإلكترون مباشرة بالتماسك العياني لمكثف أزواج كوبر. تعطي نظرية BCS فجوة الموصلية الفائقة \(\Delta\) وطول التماسك \(\xi = \hbar v_F / (\pi \Delta)\) (لموصل فائق نقي). تحدد النسبة \(\xi / a\)، حيث \(a\) هو متوسط المسافة بين الإلكترونات، عدد الإلكترونات المترابطة بشكل متماسك. عدد أزواج كوبر داخل حجم التماسك هو \(N_{\text{coh}} = (\xi/a)^3\).
	
	كسب الطاقة بسبب التكاثف هو \(\Delta E = \frac{1}{2} N(0) \Delta^2\) لكل وحدة حجم، حيث \(N(0)\) هي كثافة الحالات عند مستوى فيرمي. يمكن التعبير عن كفاءة التناسق الإضافية كـ
	\[
	K_{\mathrm{eff,SC}} - 1 = \eta \frac{\Delta E}{E_F} \cdot N_{\text{coh}}^{\beta},
	\]
	حيث \(E_F\) هي طاقة فيرمي، \(\eta\) ثابت بلا أبعاد من رتبة الواحد، و \(\beta = 2/3\) هو أس YUCT الشامل. باستخدام علاقة BCS \(\Delta = 1.76 \, k_B T_c\) وبارامترات نموذجية لـ YBCO (\(T_c \approx 92\ \text{K}\)، \(\xi \approx 2\ \text{nm}\)، \(a \approx 0.4\ \text{nm}\))، نحصل على \(\Delta E/E_F \sim 10^{-3}\) و \(N_{\text{coh}} = 125\). وبالتالي
	\[
	K_{\mathrm{eff,SC}} - 1 \sim 10^{-3} \times 125^{2/3} = 0.025,
	\]
	وهو متسق مع التقدير الظاهري في القسم~\ref{subsec:numerical}. يوضح هذا الاشتقاق أن \(K_{\mathrm{eff,mat}}\) ليس معاملًا مرتجلاً بل يتبع مباشرة من البارامترات المجهرية للموصل الفائق (الفجوة، طول التماسك، طاقة فيرمي) والأس الشامل \(\beta = 2/3\).
	
	\subsection{تجربة معملية مقترحة: قرص دوار فائق التوصيل كاختبار للهندسة المعتمدة على المادة والانهيار الحرج}
	
	يتنبأ إطار YUCT بأن الهندسة الفعالة لقرص دوار وسرعة انهياره الحرجة لا تعتمدان فقط على السرعة الزاوية ولكن أيضًا على كفاءة التناسق الداخلية للمادة \(K_{\mathrm{eff,mat}}\). هذا الاعتماد غائب في كل من الميكانيكا الكلاسيكية والنسبية العامة، مما يوفر اختبارًا نظيفًا وقابلاً للتكذيب. أدناه نفصل تجربة ملموسة ومنخفضة التكلفة باستخدام قرص من موصل فائق عند درجة حرارة عالية (YBCO).
	
	\subsubsection{الهدف التجريبي}
	
	اختبار تنبؤ YUCT بأن السرعة الزاوية الحرجة \(\omega_{\text{crit}}\) والمحيط الفعال \(C(\omega)\) يختلفان بين الحالة الطبيعية (\(T > T_c\)) والحالة فائقة التوصيل (\(T < T_c\)) لنفس القرص، بينما تُبقى جميع البارامترات الأخرى (الهندسة، الكتلة، السرعة الزاوية) متطابقة.
	
	\subsubsection{الأجهزة}
	
	\begin{itemize}
		\item قرصان متطابقان من YBCO (نصف القطر \(R = 0.1\ \text{m}\)، السمك \(1\ \text{mm}\)، مصقولان للتسطيح البصري). أحدهما كاحتياطي/تحكم.
		\item مبرد (cryostat) قادر على الحفاظ على \(T = 90\ \text{K}\) (أعلى من \(T_c \approx 92\ \text{K}\) لـ YBCO) و \(T = 77\ \text{K}\) (النيتروجين السائل، أقل من \(T_c\)).
		\item منصة دوران عالية الدقة بسرعة زاوية \(\omega\) متغيرة من 0 إلى \(2000\ \text{rad/s}\) (حوالي 19 000 دورة/دقيقة)، بدقة \(\pm 0.1\ \text{rad/s}\).
		\item مقياس تداخل ليزري (He‑Ne، \(\lambda = 632.8\ \text{nm}\)) مرتب لقياس المحيط عبر تأثير ساجناك أو بعدد الأهداب على مقياس عاكس مثبت على الحافة (دقة \(\Delta C/C \approx 10^{-6}\)).
		\item كاشفات ضوئية (ثنائيات PIN سريعة، عرض نطاق \(> 1\ \text{MHz}\)) موضوعة حول القرص لتسجيل أي انبعاث ضوئي أثناء الانهيار.
		\item كاميرا عالية السرعة (معدل إطارات \(> 10^4\ \text{إطار/ثانية}\)) لتسجيل عملية التفتت.
		\item مقاييس تسارع أو مجسات صوتية لاكتشاف لحظة الفشل بالضبط.
	\end{itemize}
	
	\subsubsection{تقديرات نسبة الإشارة إلى الضوضاء وزمن التكامل}
	
	الكميات القابلة للقياس الرئيسية هي:
	\begin{enumerate}
		\item \textbf{التغير في المحيط} \(\Delta C/C\) عند \(\omega\) ثابت بين الحالة الطبيعية وفائقة التوصيل. المقدار المتوقع: \(10^{-4}\) إلى \(10^{-2}\). لقرص نصف قطره \(R = 0.1\ \text{m}\)، \(\Delta C \sim 10^{-5}\) إلى \(10^{-3}\ \text{m}\). يستطيع مقياس تداخل ليزري مع \(\lambda = 632.8\ \text{nm}\) تمييز \(\Delta C \sim \lambda/10 \approx 6\times10^{-8}\ \text{m}\) بعد المتوسط. وبالتالي فإن الإشارة أعلى بكثير من ضوضاء الكاشف.
		
		\item \textbf{السرعة الزاوية الحرجة} \(\omega_{\text{crit}}\). الفرق المتوقع بين الحالتين الطبيعية وفائقة التوصيل هو من رتبة \(1\%\) إلى \(10\%\). يمكن لمنصة الدوران التحكم في \(\omega\) بدقة \(0.1\ \text{rad/s}\)، بينما \(\omega_{\text{crit}}\) حوالي \(10^3\ \text{rad/s}\)؛ وبالتالي الدقة النسبية \(10^{-4}\). العامل المحدد هو قابلية تكرار حدث الفشل (التشتت الإحصائي). بتكرار القياس 10–20 مرة، يمكن تخفيض الخطأ المعياري إلى أقل من \(0.5\%\).
		
		\item \textbf{انبعاث الضوء عند الانهيار}. الطاقة الإشعاعية الكلية تقدر بـ \(E_{\text{rad}} \sim \frac{1}{2} I \omega^2 \cdot f\)، حيث \(f\) هو كسر الطاقة الدورانية المحولة إلى إشعاع متماسك. حتى مع \(f = 10^{-6}\)، لقرص عزم قصوره \(I = \frac{1}{2} M R^2\) (\(M \approx 0.3\ \text{kg}\))، \(E_{\text{rad}} \sim 10^{-4}\ \text{J}\). إذا انبعث في نبضة ضيقة مدتها \(\tau \sim 1\ \mu\text{s}\)، فإن القدرة الذروية \(\sim 100\ \text{W}\). يمكن للثنائي الضوئي السريع بحساسية \(\sim 1\ \text{A/W}\) ومضخم ترانزيمبيدانس اكتشاف مثل هذه النبضة بسهولة. التحدي الرئيسي هو تمييز الانبعاث المتماسك عن إشعاع الجسم الأسود الحراري. تتنبأ YUCT بخطوط ضيقة عند ترددات فونون محددة (مثل الفونونات الضوئية عند \(\sim 10\ \text{THz}\))، والتي يمكن حلها باستخدام مطياف شبكة أو مجموعة من المرشحات ضيقة النطاق.
	\end{enumerate}
	
	\paragraph{زمن التكامل.}
	قياس المحيط يتطلب بضع ثوانٍ فقط لكل نقطة سرعة زاوية؛ قياس السرعة الحرجة يتطلب حتى 10 دقائق لكل قرص (رفع السرعة ببطء). الانبعاث الضوئي هو حدث واحد. الوقت الإجمالي للقياس لمجموعة كاملة (طبيعي + فائق التوصيل، 5–10 أقراص) أقل من أسبوع من العمل المخبري النشط.
	
	\paragraph{مصادر الضوضاء والتخفيف.}
	\begin{itemize}
		\item \textbf{الاهتزازات الميكانيكية}: استخدام طاولة بصرية معزولة عن الاهتزازات ودوار متوازن.
		\item \textbf{التمدد الحراري}: يُتحكم فيه بواسطة المبرد والقياس التفاضلي.
		\item \textbf{التداخل الكهرومغناطيسي}: كبلات محمية وقفاز فاراداي.
		\item \textbf{الضوء الخلفي}: إجراء التجارب في غرفة مظلمة؛ استخدام كشف تزامني.
	\end{itemize}
	جميع مصادر الضوضاء قابلة للإدارة بممارسات مخبرية قياسية.
	
	\subsubsection{الإجراء}
	
	\begin{enumerate}
		\item \textbf{معايرة في السكون}. قياس المحيط \(C_0(T)\) ونصف القطر \(R(T)\) للقرص عند كل من \(T = 90\ \text{K}\) و \(T = 77\ \text{K}\) مع \(\omega = 0\). هذا يحسب التمدد الحراري.
		
		\item \textbf{خط الأساس للحالة الطبيعية (\(T = 90\ \text{K}\))}. 
		\begin{enumerate}
			\item تدوير القرص بسرعة زاوية متزايدة تدريجيًا.
			\item تسجيل \(C(\omega)\) باستخدام مقياس التداخل.
			\item الاستمرار حتى يتفكك القرص. تسجيل السرعة الزاوية الحرجة \(\omega_{\text{crit, norm}}\) ونمط الفشل.
			\item إذا حدث انبعاث ضوئي، تسجيل طيفه وشدته.
		\end{enumerate}
		
		\item \textbf{الحالة فائقة التوصيل (\(T = 77\ \text{K}\))}. تكرار نفس الإجراء على قرص متطابق (أو نفس القرص بعد التسخين وإعادة التبريد).
		
		\item \textbf{التحليل التفاضلي}. مقارنة:
		\[
		\Delta \omega_{\text{crit}} = \omega_{\text{crit, SC}} - \omega_{\text{crit, norm}},\qquad
		\frac{\Delta C}{C} = \frac{C_{\text{SC}}(\omega) - C_{\text{norm}}(\omega)}{C_0(77\ \text{K})}.
		\]
	\end{enumerate}
	
	\subsubsection{تنبؤات YUCT المتوقعة}
	
	\begin{enumerate}
		\item \textbf{تغير السرعة الزاوية الحرجة}. وفقًا لصيغة YUCT
		\[
		\omega_{\text{crit}} = \frac{c}{R}\sqrt{1 - \left(\frac{K_{\mathrm{eff,mat}}}{K_{\mathrm{crit}}}\right)^2},
		\]
		يجب أن تظهر الحالة فائقة التوصيل (\(K_{\mathrm{eff,mat}}\) أعلى) \(\omega_{\text{crit}}\) مختلفًا. بالنسبة لـ YBCO، باستخدام تقديرات القسم~\ref{subsec:numerical}، التغير النسبي متوقع من رتبة \(10^{-2}\) إلى \(10^{-1}\) (بضعة بالمائة إلى عشرات بالمائة) — وهو ضمن قابلية الكشف التجريبي.
		
		\item \textbf{محيط معدل}. عند نفس \(\omega\)، يجب أن يكون المحيط الفعال في الحالة فائقة التوصيل أكبر:
		\[
		\frac{C_{\text{SC}}(\omega)}{C_{\text{norm}}(\omega)} = \frac{K_{\mathrm{eff,SR}}(\omega) K_{\mathrm{eff,mat,SC}}}{K_{\mathrm{eff,SR}}(\omega) K_{\mathrm{eff,mat,norm}}} = \frac{K_{\mathrm{eff,mat,SC}}}{K_{\mathrm{eff,mat,norm}}} > 1.
		\]
		
		\item \textbf{انبعاث ضوئي عند الانهيار}. عندما يفشل القرص، يجب أن يؤدي الفقدان المفاجئ للتناسق إلى إطلاق طاقة كإشعاع متماسك (ليس مجرد جسم أسود). تتنبأ YUCT بخطوط انبعاث ضيقة تتوافق مع رنينات الشبكة البلورية (مثل ترددات الفونونات الضوئية)، بدلاً من طيف حراري عريض. هذا توقيع فريد غائب في الكسر الكلاسيكي.
		
		\item \textbf{اختفاء المحور}. يجب أن يُظهر التصوير عالي السرعة أنه بعد الانهيار، تتحرك الشظايا دون أي مركز دوران مشترك — يتوقف المحور عن كونه خطًا مميزًا فيزيائيًا. في المقابل، فإن الجسم الجاسئ الكلاسيكي يحافظ على مركز كتلته يتحرك في خط مستقيم، لكن المحور كخط عبر ذلك المركز قابل دائمًا للتعريف. تدعي YUCT أنه بعد فقدان التناسق، لا يكون لهذا الخط معنى فيزيائي؛ يمكن اختبار ذلك بتتبع مسارات الشظايا.
	\end{enumerate}
	
	\subsubsection{مقارنة مع الفيزياء القياسية}
	
	\begin{table}[h]
		\centering
		\caption{تباين بين التنبؤات القياسية و YUCT لتجربة القرص الدوار.}
		\label{tab:experiment_predictions}
		\begin{tabular}{p{5cm}p{4cm}p{4cm}}
			\toprule
			\textbf{الملاحظة} & \textbf{الفيزياء القياسية (النسبية العامة + المرونة)} & \textbf{YUCT} \\
			\midrule
			\(\omega_{\text{crit}}\) في الحالة فائقة التوصيل مقابل العادية & متطابقة (تحددها قوة المادة فقط؛ الموصلية الفائقة لا تغير قوة الشد بشكل كبير) & مختلفة (بسبب تغير \(K_{\mathrm{eff,mat}}\)) \\
			المحيط عند \(\omega\) ثابت & متطابق (الهندسة تعتمد فقط على \(\omega\) عبر عامل لورنتز) & مختلف (عامل إضافي \(K_{\mathrm{eff,mat}}\)) \\
			انبعاث ضوئي عند الانهيار & حراري (جسم أسود) أو لا شيء & متماسك، خطوط ضيقة (رنينات بلورية) \\
			المحور بعد الانهيار & قابل دائمًا للتعريف (خط رياضي عبر مركز الكتلة) & يتوقف عن الوجود كخط مميز فيزيائيًا \\
			\bottomrule
		\end{tabular}
	\end{table}
	
	\subsubsection{الجدوى وتقدير التكلفة}
	
	\begin{itemize}
		\item \textbf{تكلفة المعدات:} المبرد (∼\$10 000)، منصة الدوران (∼\$20 000)، مقياس تداخل ليزري (∼\$15 000)، الكاشفات (∼\$5 000). المجموع ∼\$50 000 – في حدود ميزانية مختبر فيزياء أو علوم المواد النموذجي.
		\item \textbf{الزمن:} شهر للإعداد، شهر للقياسات، شهر للتحليل.
		\item \textbf{المخاطرة:} منخفضة – جميع المكونات متاحة تجاريًا، وأقراص YBCO منتجات قياسية.
	\end{itemize}
	
	\subsubsection{مسار نحو التحقق التجريبي: فرص التعاون}
	
	التجربة المقترحة تقع ضمن قدرات المختبرات الحالية ذات درجات الحرارة المنخفضة والدوران عالي السرعة. نقترح الاتصال بمجموعات بحثية لديها بالفعل خبرة في:
	\begin{itemize}
		\item \textbf{المبردات الدوارة والمواد فائقة التوصيل}: مثل مجموعة البروفيسور جون ر. هل في جامعة واشنطن (سابقًا في بوينغ)، أو مجموعة الموصلية الفائقة في معهد فريتز هابر في برلين.
		\item \textbf{منصات الدوران عالية الدقة وقياس التداخل}: مثل مجموعة الموجات الثقالية في معهد ماساتشوستس للتكنولوجيا (LIGO)، التي لديها خبرة واسعة في قياس التداخل بالليزر وعزل الاهتزازات.
		\item \textbf{التصوير عالي السرعة ودراسات التفتت}: مثل مجموعة المواد الديناميكية في معهد كاليفورنيا للتكنولوجيا أو معهد فيزياء الصدمات في جامعة ولاية واشنطن.
		\item \textbf{المؤسسات الروسية}: معهد كابيتسا للمشاكل الفيزيائية (موسكو) لديه تقليد طويل في فيزياء درجات الحرارة المنخفضة وتجارب الدوران عالي السرعة؛ معهد فيزياء الحالة الصلبة (تشيرنوجولوفكا) يعمل بشكل مكثف مع YBCO وغيره من الموصلات الفائقة ذات \(T_c\) العالية.
	\end{itemize}
	
	نحن منفتحون على ترتيبات تعاونية: يمكن لـ YUCT تقديم دعم نظري مفصل، بما في ذلك تنبؤات دقيقة لـ \(\omega_{\text{crit}}\) لهندسات ومواد محددة للأقراص، بالإضافة إلى المساعدة في تفسير البيانات. النتيجة الإيجابية ستكون اكتشافًا تاريخيًا: أول عرض لهندسة الزمكان المعتمدة على المادة وأول رصد لانتقال طوري منسق في نظام دوار. ندعو الباحثين التجريبيين المهتمين للاتصال بالمؤلف على \texttt{alexey@yakushev.eu}.
	
	\subsubsection{خلاصة التجربة المقترحة}
	
	إذا لوحظت الاختلافات المتوقعة بين الحالتين الطبيعية وفائقة التوصيل، فإن ذلك سيوفر تأكيدًا مخبريًا مباشرًا لادعاء YUCT أن الهندسة معتمدة على المادة وأن محور الدوران هو خاصية ناشئة للمادة المنسقة. النتيجة الصفرية (لا فرق) ستضع حدًا أعلى لـ \(K_{\mathrm{eff,mat}} - 1\) لـ YBCO، مما يقيد النظرية. التجربة بسيطة وبأسعار معقولة وحاسمة – إنها تحول مفارقة إرنفست من تجربة فكرية إلى منصة اختبار تجريبية لنموذج التناسق.
	
	\subsection{عواقب أنطولوجية: المسافة، الزمن، درجة الحرارة}
	
	يجبرنا نموذج التناسق على إعادة النظر في مكانة الكميات الفيزيائية الأكثر أساسية.
	
	\begin{itemize}
		\item \textbf{المسافة.} في YUCT، تُعرَّف المسافة الفعالة بين مجموعتي تناسق \(i\) و \(j\) على أنها \(d_{ij} = c \cdot \langle \tau_{ij} \rangle\)، حيث \(\tau_{ij}\) هو متوسط الوقت اللازم لدليل قصير للانتشار بين المجموعتين. ينبثق موتر المترية \(g_{\mu\nu}(x)\) بعد ترفيع (compactification) حقل التناسق \(\Psi_{MN}\) ذي الـ 19 بُعدًا ثم متوسط الدرجات الداخلية. وبالتالي \textbf{المسافة ليست معطاة أولية بل مقياس إحصائي لعمليات تبادل الأدلة}. بدون مادة منسقة لا يوجد معنى فيزيائي لكلمة ``مسافة''.
		
		\item \textbf{الزمن.} كما هو موضح في الملحق الخامس، الزمن الذاتي يتناسب مع زيادة إنتروبيا التناسق: \(d\tau \propto dS_{\mathrm{coord}}\). الوكيل المعزول الذي لا يشارك في أي عملية تناسق لا يختبر الزمن. \textbf{الزمن هو مقياس لنشاط التناسق، وليس معاملًا مطلقًا خارجيًا}.
		
		\item \textbf{درجة الحرارة.} درجة الحرارة هي إسقاط عياني لمتوسط خطأ التناسق \(\varepsilon\). قانون الخطأ الشامل \(\varepsilon = \kappa_c \alpha K_{\mathrm{eff}}^{-\beta}\) مع نظرية التناسق الأساسية (الملحق ب) يضمن \(\varepsilon > 0\) لأي \(K_{\mathrm{eff}}\) محدود. وبالتالي \textbf{الصفر المطلق لدرجة الحرارة غير قابل للتحقيق}، ليس بسبب قيد تقني بل لأن الحالة المنسقة تمامًا مستحيلة من حيث المبدأ.
	\end{itemize}
	
	باختصار، المسافة والزمن ودرجة الحرارة ليست لبنات بناء أولية للواقع. إنها \textbf{بارامترات إحصائية ناشئة} لشبكة تناسق، تظهر فقط عندما تنسق المادة حالاتها عبر تبادل الأدلة وتفعيل القواميس المشتركة.
	
	\section{اعتراضات وردود}
	\label{sec:objections}
	
	\paragraph{الاعتراض 1: ``مجرد إعادة تسمية لـ \(\gamma\)''}
	النسبية العامة تستخدم بالفعل \(\gamma\) لوصف القرص. يبدو أن YUCT تستبدل \(\gamma\) بـ \(K_{\mathrm{eff}}\) دون محتوى جديد.
	
	\paragraph{الرد}
	في النسبية العامة، \(\gamma\) محدد بشكل فريد بواسطة \(\omega\) وسرعة الضوء في الفراغ. تقدم YUCT عاملًا إضافيًا \(K_{\mathrm{eff,mat}}\) يمكن تغييره بشكل مستقل (على سبيل المثال، بواسطة التماسك الكمي). هذا يؤدي إلى انحراف قابل للاختبار عن النسبية العامة: نفس \(\omega\) يعطي محيطات مختلفة اعتمادًا على كفاءة تناسق المادة. علاوة على ذلك، فإن التعريف \(K_{\mathrm{eff}} = \gamma\) ليس مجرد إعادة تسمية؛ إنها صيغة محددة تمليها ثوابت التناسق الشاملة \(\Sodd\) و \(\Seven\) والتقارب \(\Sodd \varphi^2 \approx \pi\). توفر هذه الثوابت مبررًا هندسيًا أعمق غائبًا في النسبية العامة.
	
	\paragraph{الاعتراض 2: ``انتهاك السببية''}
	كيف يمكن لجميع نقاط القرص معرفة القاموس بدون إشارات فائقة الضوء؟
	
	\paragraph{الرد}
	يُوزع القاموس مسبقًا عبر البنية المادية (مثلًا، ثوابت الشبكة، ترتيب السبين). عندما يبدأ الدوران، تفرض التآثرات الكهرومغناطيسية والمرنة المحلية المسافات المقررة. هذا ليس أكثر من إشارات فائقة الضوء من الجسم الجاسئ في الميكانيكا النيوتونية — القيود مبنية مسبقًا، وليس يتم التواصل بها في الوقت الفعلي.
	
	\paragraph{الاعتراض 3: ``متغيرات خفية محظورة بنظرية بل''}
	يشبه القاموس الموزع مسبقًا متغيرًا خفيًا تستبعده ميكانيكا الكم للتشابك.
	
	\paragraph{الرد}
	القاموس في YUCT ليس متغيرًا خفيًا محليًا بمعنى نظرية بل. إنه بنية ارتباطية عالمية غير قادرة على الإشارات، تُنشأ خارج الخط ولا يمكن استخدامها لنقل المعلومات أسرع من الضوء. في ميكانيكا الكم، الحالة المتشابكة نفسها هي مثل هذه الارتباطية. YUCT ببساطة توسع هذا المفهوم إلى مجال التناسق الكلاسيكي والنسبوي. تمنع نظرية بل المتغيرات الخفية الواقعية المحلية، لكن YUCT لا تفترض الواقعية المحلية؛ إنها تفترض التناسق كأولية، والقاموس صريحًا غير محلي بمعنى أنه موزع مسبقًا عبر النظام. لا يمكن الإشارة فائقة الضوء لأن الدليل (الدوران نفسه) لا يزال ينتشر بشكل سببي. وبالتالي، فإن YUCT متوافقة تمامًا مع جميع الاختبارات التجريبية لمتباينات بل.
	
	\paragraph{الاعتراض 4: ``لا توجد آلية فيزيائية لتغيير \(K_{\mathrm{eff,mat}}\)''}
	كيف يمكن تغيير كفاءة التناسق دون تغيير \(\omega\)؟
	
	\paragraph{الرد}
	تشمل الأمثلة: التماسك فائق التوصيل (حالة كمومية عيانية)، والتذبذبات الكمية المستحثة بالحقل المغناطيسي، والمشابك البصرية. هذه تؤثر مباشرة على قدرة المادة على الحفاظ على مسافات متبادلة دقيقة، أي كفاءة تناسقها. أس القياس الشامل \(\betaf = 2/3\) من الملحق ل يقدم رابطًا كميًا بين حجم التماسك والتغير المتوقع في \(K_{\mathrm{eff,mat}}\).
	
	\paragraph{الاعتراض 5: ``صيغة السرعة الحرجة تختلف عن ميكانيكا الكسر الكلاسيكية''}
	تتنبأ YUCT بسرعة حرجة معتمدة على المادة يمكن أن تنحرف عن الصيغة الكلاسيكية. كيف يمكن التوفيق بين هذا والمعرفة الهندسية الراسخة؟
	
	\paragraph{الرد}
	بالنسبة للمواد العادية (غير فائقة التوصيل، غير متماسكة)، \(K_{\mathrm{eff,mat}} = 1\) وتنخفض صيغة YUCT إلى النتيجة الكلاسيكية تمامًا عندما يُعبر عن \(\Kcrit\) بدلالة قوة الشد. الانحراف يظهر فقط عندما \(K_{\mathrm{eff,mat}} \neq 1\). وبالتالي فإن YUCT لا تتعارض مع الميكانيكا الكلاسيكية؛ بل توسعها إلى أنظمة جديدة. النظام الكلاسيكي هو حالة خاصة.
	
	\section{ملخص بياني: مخطط الطور لـ \(K_{\mathrm{eff}}\)}
	\label{sec:diagram}
	
	\begin{figure}[h]
		\centering
		\begin{tikzpicture}[scale=1.2]
			% المحاور
			\draw[->] (0,0) -- (6,0) node[right] {\(K_{\mathrm{eff}}\)};
			\draw[->] (0,0) -- (0,4) node[above] {الهندسة / الفيزياء};
			% المناطق
			\fill[blue!10] (0,0) rectangle (2,4);
			\fill[green!15] (2,0) rectangle (4.5,4);
			\fill[red!10] (4.5,0) rectangle (6,4);
			% التسميات
			\node at (1,3.5) {كلاسيكي};
			\node at (1,3) {\(K_{\mathrm{eff}} = 1\)};
			\node at (1,2.5) {إقليدي};
			\node at (1,2) {زمن مطلق};
			\node at (3.25,3.5) {نسبوي};
			\node at (3.25,3) {\(K_{\mathrm{eff}} > 1\)};
			\node at (3.25,2.5) {غير إقليدي};
			\node at (3.25,2) {فضاء منحنٍ};
			\node at (5.25,3.5) {كمومي};
			\node at (5.25,3) {\(K_{\mathrm{eff}} \to \infty\)};
			\node at (5.25,2.5) {غير تبادلي};
			\node at (5.25,2) {تشابك};
			% خط العتبة الحرجة
			\draw[dashed, thick, orange] (2.5,0) -- (2.5,4);
			\node[orange, rotate=90] at (2.5,2) {\(K_{\mathrm{crit}}\)};
			% أسهم تشير إلى تطور القرص
			\draw[thick, ->] (0.5,0.5) to[bend left] (5.5,0.5);
			\node at (3,0.2) {زيادة \(\omega\) أو \(K_{\mathrm{eff,mat}}\)};
			% خطوط رأسية عند Keff=1 و Keff=∞
			\draw[dashed] (2,0) -- (2,4);
			\draw[dashed] (4.5,0) -- (4.5,4);
			\node at (2, -0.3) {1};
			\node at (4.5, -0.3) {\(\infty\)};
		\end{tikzpicture}
		\caption{مخطط الطور للقرص الدوار في YUCT. نفس المعامل \(K_{\mathrm{eff}}\) يحكم الانتقال من الكلاسيكي (\(K_{\mathrm{eff}}=1\)) عبر النسبوي (\(K_{\mathrm{eff}}>1\)) إلى الكمومي (\(K_{\mathrm{eff}}\to\infty\)). العتبة الحرجة \(K_{\mathrm{crit}}\) تمثل الحد الفاصل بين التناسق المستقر (يسار) وفقدان التناسق الكارثي (يمين). بالنسبة لـ \(K_{\mathrm{eff,total}} < K_{\mathrm{crit}}\) (ما وراء الخط البرتقالي)، يتفكك القرص.}
		\label{fig:phasediagram}
	\end{figure}
	
	\section{الاتصال بالدوران الكوني (الملحق~\(\Omega\)) والأجسام الفيزيائية الفلكية}
	\label{sec:cosmic}
	
	تنطبق نفس مبادئ التناسق على دوران الكون ككل (انظر الملحق~\(\Omega\)). يمكن تفسير ثنائي القطب في الخلفية الميكروية الكونية (CMB) كتراكب للحركة المحلية والدوران العالمي، مع كفاءة تناسق \(K_{\mathrm{eff,cosmic}}\) متعلقة بالسرعة الزاوية \(\omega_{\text{univ}}\). يوفر التحليل الحالي للقرص الدوار نظيرًا محليًا: المترية الفعالة لكون دوار ستصفها بالمثل \(K_{\mathrm{eff}}(r)\) المعتمدة على الموضع. وبالتالي فإن التنبؤ القابل للاختبار بهندسة معتمدة على المادة لقرص دوار له نظير كوني: قد يعتمد اتساع ثنائي القطب لـ CMB على كفاءة التناسق على نطاق واسع للكون. على الرغم من عدم قابليته للاختبار المباشر في المختبر، إلا أن هذا الرابط المفاهيمي يعزز الاتساق الداخلي لـ YUCT.
	
	النجوم النيوترونية (النجوم النابضة، المغناطيسارات) هي نظائر فيزيائية فلكية طبيعية. دورانها السريع (حتى \(\sim 700\) هرتز) وحقولها المغناطيسية القوية تخلق بيئة قصوى حيث تُحدد كفاءة التناسق بواسطة الجاذبية، والقوى النووية، وحفظ التدفق المغناطيسي. يمكن تطبيق شرط الانهيار الحرج المشتق في القسم~\ref{sec:critical-disruption} لتقدير أقصى تردد دوران للنجم النيوتروني. باستخدام صيغة YUCT مع بارامترات مادة مناسبة (الكثافة \(\sim 10^{17}\) كغم/م³، قوة الشد \(\sim 10^{30}\) باسكال، \(K_{\mathrm{eff,mat}}\) ربما معززة بفعل السيولة الفائقة)، نحصل على تردد أقصى حوالي 1 كيلوهرتز، وهو متوافق مع الملاحظات. وهذا يوفر تدقيقًا متقاطعًا مستقلاً للإطار.
	
	\subsection{الاتصال بالحلقة الجبرية لـ YUCT: لماذا الهندسة معتمدة على المادة}
	\label{subsec:ehrenfest-algebraic-loop}
	
	كشف تحليل القرص الدوار أن الهندسة المكانية الفعالة — وتحديدًا نسبة المحيط إلى نصف القطر — ليست خاصية مطلقة للزمكان بل تعتمد على كفاءة التناسق \(K_{\mathrm{eff,mat}}\) للمادة. يبرز سؤال طبيعي: هل \(K_{\mathrm{eff,mat}}\) معامل قابل للتعديل بحرية، أم أنه مقيد بالبنية الجبرية الأعمق لـ YUCT؟
	
	يكمن الجواب في الثوابت الأساسية المشتقة في الملحق Ferra1.2. التحلل الهرمي للتناسق يعطي ثابتين عالميين:
	\[
	S_{\mathrm{odd}} = \frac{6}{5} = 1.2,\qquad S_{\mathrm{even}} = \frac{4}{5} = 0.8,
	\]
	اللذان يحققان العلاقات الدورية \(S_{\mathrm{odd}}+S_{\mathrm{even}}=2\) و \(S_{\mathrm{odd}}/S_{\mathrm{even}} = 3/2 = 1/\beta\). تحكم هذه الثوابت تكميم نصف الصحيح لكتل الفرميونات (\(m_f \propto q^{N_f}\) مع \(q = (3/2)^{1/3}\)) كما تملي أيضًا سلوك التناسق الحدي في الأنظمة العيانية.
	
	من اللافت أن نفس البنية الجبرية توفر تقريبًا عالي الدقة لثابت الدائرة \(\pi\). باستخدام النسبة الذهبية \(\varphi = (1+\sqrt{5})/2\)، نجد
	\begin{equation}
		S_{\mathrm{odd}} \cdot \varphi^2 = \frac{6}{5} \left( \frac{1+\sqrt{5}}{2} \right)^2 
		= \frac{9+3\sqrt{5}}{5} \approx 3.1416407865\dots
	\end{equation}
	وهذا يختلف عن \(\pi \approx 3.1415926535\) بنسبة \(4.8\times10^{-5}\) فقط (خطأ نسبي \(1.5\times10^{-5}\)). وهذا أكثر دقة من التقريب الكسري الكلاسيكي \(22/7\) بمرتبة مقدارية.
	
	ضمن YUCT، هذا التقارب ليس مصادفة. النسبة الهندسية الفعالة \(\pi_{\mathrm{eff}}\) لنظام ذي كفاءة تناسق محدودة \(K_{\mathrm{eff}}\) تُعطى بواسطة
	\begin{equation}
		\pi_{\mathrm{eff}}(K_{\mathrm{eff}}) = \pi_{\infty} - \Delta\pi \cdot K_{\mathrm{eff}}^{-\beta},
	\end{equation}
	حيث \(\pi_{\infty} = \pi\) هي نهاية التناسق التام (\(K_{\mathrm{eff}} \to \infty\))، و \(\Delta\pi\) ثابت يتعلق بالبنية المنفصلة للقاموس. بالنسبة لمادة في حالتها الأرضية، يُحدد \(K_{\mathrm{eff}}\) بالمستوى الهرمي السائد، الذي يتميز بالنسبة \(S_{\mathrm{odd}}/S_{\mathrm{even}} = 3/2\). بتعويض هذا في قانون القياس نحصل على \(\pi_{\mathrm{eff}} \approx S_{\mathrm{odd}}\varphi^2\).
	
	\textbf{الآثار المترتبة على القرص الدوار.}
	كفاءة التناسق المادية \(K_{\mathrm{eff,mat}}\) التي تعدل المحيط في المعادلة~(\ref{eq:keff-disk}) ليست إذن معامل ملاءمة عشوائيًا بل هي متجذرة في نفس الحلقة الجبرية التي تحدد كتل الجسيمات الأولية. الانحراف المتوقع لـ \(\pi\) الفعال عن قيمته الرياضية صغير جدًا في الظروف العادية (بما أن \(K_{\mathrm{eff,mat}}\) كبير)، وهذا يفسر لماذا الهندسة الإقليدية تقريب ممتاز للهندسة اليومية. ومع ذلك، في الأنظمة القصوى — مثل قرب انتقال طوري حيث ينخفض \(K_{\mathrm{eff,mat}}\) بشكل حاد — قد يصبح الانحراف قابلًا للكشف. وهذا يوفر رابطًا مباشرًا بين مفارقة إرنفست، ومسألة التسلسل الهرمي لفيزياء الجسيمات، والتجربة المخبرية المقترحة بقرص فائق التوصيل (القسم~\ref{sec:experiment}).
	
	باختصار، الاعتماد المادي للهندسة ليس افتراضًا مرتجلًا لـ YUCT بل نتيجة ضرورية لنفس البنية الجبرية المنفصلة التي توحد طيف الفرميونات الأساسية ونشوء الزمكان المستمر.
	
	\section{الاستنتاج}
	\label{sec:conclusion}
	
	قدمنا تفسيرًا قائمًا على التناسق لمفارقة إرنفست ضمن إطار YUCT، المرتكز الآن على ثوابت التناسق الشاملة \(\Sodd = 1.2\)، \(\Seven = 0.8\)، والتقارب \(\Sodd \varphi^2 \approx \pi\). يُنظر إلى القرص الدوار كنظام يتشارك قاموسًا موزعًا مسبقًا للمسافات المتبادلة؛ يعمل الدوران كدليل يُفعِّل مدخلات القاموس، مؤديًا إلى مترية مكانية فعالة غير إقليدية. تعطي كفاءة التناسق \(K_{\mathrm{eff}}(r)\) مباشرة العامل الذي ينحرف به المحيط عن القيمة الإقليدية. عندما \(K_{\mathrm{eff}} \to 1\)، تصبح الهندسة إقليدية والزمن مطلقًا، مستعيدًا الفيزياء النيوتونية.
	
	قدمنا مفهومي \textbf{مركز التناسق} (محور الدوران) كخاصية ناشئة للمادة المنسقة، و \textbf{الانهيار الحرج} كفقدان كارثي للتناسق عندما تنخفض \(\Kefftot\) عن عتبة معتمدة على المادة \(\Kcrit\). اشتقنا صيغة كمية للسرعة الزاوية الحرجة وأظهرنا أنها تنخفض إلى ميكانيكا الكسر الكلاسيكية عندما \(\Keffmat=1\). يتنبأ الإطار بأن المواد ذات التناسق الداخلي المعزز (مثل الموصلات الفائقة) يجب أن تظهر سرعة حرجة مختلفة وهندسة فعالة معدلة — فرق قابل للاختبار عن كل من الفيزياء الكلاسيكية والنسبية العامة.
	
	لا تتعارض YUCT مع النسبية العامة؛ بل تقدم تفسيرًا أعمق لأصل الانحناء كخاصية ناشئة للتناسق عالي الكفاءة، مع توفر الثوابت الهرمية للأساس الكمي. يضمن التقارب \(\Sodd \varphi^2 \approx \pi\) (خطأ نسبي \(1.5\times10^{-5}\)) أن الهندسة الإقليدية تُستعاد بدقة عالية عندما يكون التناسق ضئيلًا، وأي انحراف تتحكم فيه نفس الثوابت الهرمية. نتيجة تجريبية إيجابية من تجربة القرص فائق التوصيل المقترحة لن تؤكد YUCT فحسب، بل ستثبت أيضًا أنه يمكن التحكم في هندسة الزمكان عبر حالات المادة الكمومية، مما يفتح طريقًا جديدًا للفيزياء الأساسية.
	
	\begin{thebibliography}{99}
		\bibitem{Ehrenfest1909}
		P. Ehrenfest, ``Gleichf{\"o}rmige Rotation starrer K{\"o}rper und Relativit{\"a}tstheorie,'' \emph{Physikalische Zeitschrift} \textbf{10}, 918 (1909).
		
		\bibitem{Yakushev2026}
		A. V. Yakushev, \emph{Yakushev Unified Coordination Theory (YUCT)}, Zenodo, \url{https://doi.org/10.5281/zenodo.18444598} (2026).
		
		\bibitem{AppendixB}
		الملحق ب: مفارقة التناسق الزمني -- YUCT.
		
		\bibitem{AppendixL}
		الملحق ل: قانون خطأ القياس الكسري -- قانون شامل من الدنا إلى علم الكون.
		
		\bibitem{AppendixR}
		الملحق ر: التحكم التناسقي في العمليات النووية -- من الاضمحلال المتحكم فيه إلى الاندماج البارد.
		
		\bibitem{AppendixG}
		الملحق ز: نظرية ياكوشيف الموحدة للتناسق -- حل أساسي لمشكلة الجاذبية الكمية.
		
		\bibitem{AppendixΩ}
		الملحق~\(\Omega\): الدوران العالمي للكون وعمره الأدنى الجديد من نصف قطر دوران ثنائي القطب.
		
		\bibitem{AppendixFerra1.2}
		الملحق Ferra1.2: ثوابت التناسق الشاملة للأطوار المنتظمة وغير المنتظمة.
		
		\bibitem{Langevin1935}
		P. Langevin, ``La notion de corps rigide en relativit{\'e},'' \emph{Comptes rendus} \textbf{200}, 1900--1903 (1935).
		
		\bibitem{Gron1995}
		{\O}. Gr{\o}n, ``The relativistic rigid rotating disk,'' \emph{Physics Essays} \textbf{8}, 92--102 (1995).
		
		\bibitem{M87Precession2024}
		Cui, Y. et al., ``Precessing jet in M87: Evidence for a 11-year period from VLBI monitoring 2000–2022,'' \emph{Astrophysical Journal} \textbf{960}, 125 (2024).
		
		\bibitem{NatureAstronomy2025}
		Bright, J. S. et al., ``Transient jet speeds in stellar-mass black holes and neutron stars: The largest sample to date,'' \emph{Nature Astronomy} \textbf{9}, 112–128 (2025).
		
		\bibitem{JetQuenching2006}
		Fender, R. P., Belloni, T. M., \& Gallo, E., ``Towards a unified model for black hole X-ray binaries,'' \emph{Monthly Notices of the Royal Astronomical Society} \textbf{363}, 1105–1118 (2006).
		
		\bibitem{JetQuenching2018}
		Tetarenko, A. J. et al., ``The jet quenching factor in black hole X-ray binaries,'' \emph{Astrophysical Journal} \textbf{862}, 90 (2018).
		
		\bibitem{TDE2016}
		Pasham, D. R. et al., ``A relativistic jet from a tidal disruption event,'' \emph{Astrophysical Journal} \textbf{820}, L18 (2016).
		
		\bibitem{TDE2020}
		Alexander, K. D. et al., ``The tidal disruption event AT2019dsg: A jetted outflow on the rise,'' \emph{Astrophysical Journal Letters} \textbf{894}, L19 (2020).
		
		\bibitem{Yarman2013}
		T. Yarman, ``Radiation from an accelerating neutral body: The case of rotation,'' \emph{European Physical Journal Plus} \textbf{128}, 134 (2013). DOI:10.1140/epjp/i2013-13134-9
		
		\bibitem{Kholmetskii2020}
		A. L. Kholmetskii, T. Yarman, O. Yarman, M. Arik, ``Analyses of Mössbauer experiments in a rotating system: Proper and improper approaches,'' \emph{Annals of Physics} \textbf{418}, 168191 (2020). DOI:10.1016/j.aop.2020.168191
		
		\bibitem{YARK2016}
		A. L. Kholmetskii, T. Yarman, O. Yarman, M. Arik, ``Pound–Rebka result within the framework of YARK theory,'' \emph{Canadian Journal of Physics} \textbf{94}, 806–810 (2016). DOI:10.1139/cjp-2016-0087
		
		\bibitem{Bashkov2000}
		V. Bashkov, M. Malakhaltsev, ``Geometry of rotating disk and the Sagnac effect,'' arXiv:gr-qc/0011061 (2000).
		
	\end{thebibliography}
	
\end{document}